\documentclass[11pt,oneside]{book}
\usepackage[margin=1.1in]{geometry}
\usepackage{amsmath,amssymb,amsthm}
\usepackage{booktabs}
\usepackage{longtable}
\usepackage[hidelinks]{hyperref}
\usepackage{url}

\theoremstyle{plain}
\newtheorem{theorem}{Theorem}[chapter]
\newtheorem{proposition}[theorem]{Proposition}
\newtheorem{lemma}[theorem]{Lemma}
\newtheorem{corollary}[theorem]{Corollary}
\newtheorem{conjecture}[theorem]{Conjecture}
\theoremstyle{definition}
\newtheorem{definition}[theorem]{Definition}
\theoremstyle{remark}
\newtheorem{remark}[theorem]{Remark}
\newtheorem{observation}[theorem]{Observation}

\newcommand{\N}{N}
\newcommand{\g}{\mathfrak{g}}
\newcommand{\Ehat}{\widehat{E}}

\title{\bf Goldbach Deserts\\[4pt]
\large Construction, Optimisation, and the Limits of Covering}
\author{Anonymous}
\date{July 2026}

\begin{document}
\frontmatter
\maketitle

\chapter*{Abstract}
\addcontentsline{toc}{chapter}{Abstract}

For an even integer $\N$, let $\g(\N)$ denote the least prime occurring in a
Goldbach partition of $\N$. Empirically $\g$ is minuscule---it never exceeds
$9781$ for any even $\N\le4\cdot10^{18}$---and the interesting question is
therefore extremal and constructive: how large can $\g(\N)$ be \emph{forced}
to be, and what must be paid in the size of $\N$? This thesis studies that
question along three record games: the unbounded \emph{height} game, the
\emph{threshold} game $T(Q)$ (minimise $\N$ subject to $\g(\N)>Q$), and the
\emph{budget} game $R(X)$ (maximise $\g(\N)$ subject to $\N<X$).

Chapters~\ref{ch:foundations} and~\ref{ch:engineering} develop the
constructive theory. A congruence cover on the prime offsets, realised
through the Chinese Remainder Theorem and completed by a progression search,
produces certified integers far smaller than the classical primorial or
complete-cover constructions. We give a cost model for such searches---a
single \emph{failure exponent} $E=|U|\cdot\mathrm{boost}/\log\N$ whose
exponential $e^{E}$ predicts the number of candidates that must be
examined---and validate it against $1.2\times10^{9}$ scanned candidates, where
the measured and predicted exponents agree to within $0.1$ nats. The
programme's certified incumbents are a $2\,480$-digit integer with
$\g(\N)=1\,157\,341$, a $150$-digit integer with $\g(\N)=104\,527$, and a
$199$-digit integer with $\g(\N)=119\,419$.

Chapters~\ref{ch:frontier}--\ref{ch:closure} ask where this stops, and this
is where the thesis makes its main contributions. We show that the difficulty
of the whole enterprise is an \emph{integrality gap}: the natural
linear-programming relaxation of the covering problem is
\emph{vacuous}---fractionally, the available moduli cover $100\%$ of the prime
offsets---so the entire failure exponent is the price of rounding a fractional
cover to an integral one. We then attack that gap with exact and global
optimisation (large-neighbourhood search with certified windows, belief
propagation with decimation, banded constructions) and find it immovable at
the scales tested, with evidence that the greedy construction already sits at
or near the quenched optimum.

Finally we prove, at the level of rigour of a well-supported research
programme rather than a finished theorem, a \emph{closure} result: within a
concrete language of certificate schemas, no construction beats covering.
Each family closes for a different and quantifiable reason---a value-set
ceiling of $\sqrt{Q}$, a supply-limited purchase plan that runs $2.5\times$
over the available design entropy, a cluster bound, a lattice-isolation
argument, or an empty intersection at record scale. The theorem is
conditional on two named hypotheses, one of which---the absence of a
compositeness test cheaper than a single modular exponentiation---we isolate
and pose as an open problem of independent interest.

The closure theorem was stress-tested by an adversarial search
(an LLM-driven evolutionary system) against a machine-checkable
scorer. Its first success was not against the theorem but against our
\emph{referee}: it found a $50\times$ accounting error in the evaluator and
exploited it before any claim was published. The audit protocol caught it,
and repairing the referee strengthened the mathematics. We record this as a
methodological finding: adversarial search against a formal claim is most
valuable as a referee-hardening device.

\tableofcontents

\mainmatter
\chapter{Introduction: three games and one exponent}\label{ch:intro}

Goldbach's conjecture asserts that every even integer $\N>2$ is a sum of two
primes; it has been verified for all $\N\le4\cdot10^{18}$~\cite{OeS}.
Heuristically one expects not just one representation but many. The
Hardy--Littlewood circle method---which estimates such counts by treating the
primes as a random set of the appropriate density and then correcting for
divisibility constraints---predicts~\cite{HL}
\begin{equation}\label{eq:HL}
r(\N)\;\sim\;2\,\Pi_2\,\frac{\N}{(\log \N)^2}\prod_{\substack{p\mid \N\\ p>2}}\frac{p-1}{p-2},
\qquad
\Pi_2=\prod_{p>2}\Bigl(1-\tfrac{1}{(p-1)^2}\Bigr),
\end{equation}
where $r(\N)=\#\{p\le\N:p\text{ and }\N-p\text{ both prime}\}$. Reading this
from left to right: $\N/(\log\N)^2$ is the naive guess, since a number near
$\N$ is prime with probability about $1/\log\N$ and there are $\N$ choices for
$p$; the twin-prime constant $\Pi_2\approx0.6601$ corrects for the fact that
primes are not truly random modulo small numbers; and the last factor inflates
the count when $\N$ is highly composite. The conclusion is that a typical large
even integer should have \emph{many} Goldbach partitions. So the interesting
question is not whether a representation exists, but how lopsided the most
lopsided one must be.

\begin{definition}
For an even integer $\N\ge4$ admitting a Goldbach representation, the
\emph{least Goldbach summand} is
\[
\g(\N)\;=\;\min\{\,p\ \text{prime}:\N-p\ \text{is prime}\,\}.
\]
\end{definition}

This is the quantity $p(n)$ of the \emph{minimal Goldbach partition} studied by
Granville, van de Lune and te Riele~\cite{GLvdtR}, who conjectured a ceiling
\begin{equation}\label{eq:GLR}
\g(\N)\;=\;O\bigl((\log \N)^2\log\log \N\bigr).
\end{equation}
Empirically $\g$ is tiny and grows very slowly; its record values are
catalogued in the OEIS~\cite{OEIS} (A020481, with records A025018/A025019).
Exhaustive search gives a slowly climbing ladder: $\g=1093$ at
$\N=60\,119\,912$; $\g(\N)\le5569$ for all even
$\N\le4\cdot10^{14}$~\cite{Richstein}; and the current exhaustive record
$\g=9781$, first attained at $\N=3\,325\,581\,707\,333\,960\,528$, with
$\g(\N)\le9781$ for \emph{all} even $\N\le4\cdot10^{18}$~\cite{OeS}. Strikingly,
these records occur at integers of at most $19$ digits.

We ask the opposite, extremal and constructive question:
\begin{quote}\itshape
How large can $\g(\N)$ be forced to be, and at what cost in the size of $\N$?
\end{quote}
An even integer $\N$ for which $\g(\N)$ is large is a \emph{Goldbach desert}:
a number all of whose small prime offsets are blocked. Constructing deserts is
the subject of this thesis, and the reason the subject is interesting is that
every construction turns out to be governed by a single number---an exponent
that measures how much luck the construction still needs after all its
cleverness is spent.

\section{The three games}

There are three natural record questions, and they reward different
constructions. The first is the unbounded \emph{height game}: exhibit as large
a value of $\g(\N)$ as possible, with no limit on the size of $\N$. An ordinary
prime gap is one way to play, but it is a strictly stronger construction: a gap
must block \emph{every} integer offset, while a desert needs to block only the
\emph{prime} offsets. Prime-gap constructions are therefore a strict subclass of
the constructions available here---a point we return to in
\S\ref{sec:gaps} with a concrete two-orders-of-magnitude example.

For comparisons in which size matters we use two finite games.

\begin{definition}[Threshold and budget games]\label{def:games}
For $Q\ge2$ and $X\ge4$, define
\[
T(Q)=\min\{\,\N \text{ even, Goldbach-representable}:\ \g(\N)>Q\,\},
\qquad
R(X)=\max\{\,\g(\N):\ \N<X\,\}.
\]
The \emph{threshold game} fixes $Q$ and tries to lower $T(Q)$; the
\emph{budget game} fixes $X$ and tries to raise $R(X)$. ``Fewer than $200$
decimal digits'' means the strict budget $\N<10^{199}$.
\end{definition}

\begin{remark}[A naming caution]
The budget game has been written both as $R(10^{199})$ (``fewer than $200$
digits'') and as $R(10^{200})$ (``at most $200$ digits'') in the literature
this thesis builds on, and the two are genuinely different: our $200$-digit
incumbent with $\g=112\,249$ satisfies $\N<10^{200}$ but exceeds $10^{199}$ by
$3.4\%$, which is precisely why a $199$-digit integer with the smaller value
$\g=119\,419$ is the record for the stricter budget. Throughout we state the
bound explicitly with each record.
\end{remark}

\section{Certified incumbents}

Table~\ref{tab:records} lists the programme's certified incumbents at the time
of writing. Every entry has been verified twice: once by the pipeline that
produced it and once by an independent verifier written against the record file
rather than against the search code, including a re-derivation of every
per-offset witness and an independent check of the elliptic-curve primality
certificate.

\begin{table}[h]
\centering\small
\caption{Certified incumbents. ``Witnesses'' counts prime offsets excluded by
an explicit congruence divisor plus those excluded by a strong base-$2$
compositeness witness; in every case the positive side is an exhibited prime
complement with an ECPP certificate.}\label{tab:records}
\begin{tabular}{lrrl}
\toprule
Game & Digits of $\N$ & $\g(\N)$ & Construction\\
\midrule
Height (unbounded) & $2\,480$ & $1\,157\,341$ & cover at $Q=1.16\cdot10^6$, $t=14\,544$\\
Height, gap-dominating & $2\,692$ & $1\,134\,871$ & cover at $Q=1\,113\,138$, $k=14$\\
Threshold $T(100\,000)$ & $150$ & $104\,527$ & cover at $Q=100\,003$, $k=6.21\cdot10^{10}$\\
Budget $R(10^{199})$ & $199$ & $119\,419$ & cover at $Q\approx1.19\cdot10^5$, $t=1.46\cdot10^8$\\
Budget $R(10^{200})$ & $200$ & $112\,249$ & cover at $Q=107\,720$, $k=56\,770$\\
\bottomrule
\end{tabular}
\end{table}

Two of these deserve comment. The $2\,692$-digit entry is included because it
settles the comparison with ordinary prime gaps in the strongest possible way:
its $\g$ exceeds the largest \emph{certified} prime gap known
($1\,113\,106$), and it does so with an integer roughly $6.9$ times shorter
than the one realising that gap. Deserts really are cheaper than gaps, by a
measurable factor.

The $150$-digit threshold record is the endpoint of a \emph{ratchet}: the same
covering machinery, re-tuned and re-run, produced a descending ladder
$195\to179\to150$ digits at $Q\approx10^5$, banking a certified record at each
rung. Chapter~\ref{ch:engineering} argues that this ratchet discipline, rather
than any single clever cover, is the strategy that actually produces records
under a fixed compute budget.

\section{The exponent that governs everything}

The technical spine of the thesis is a single quantity. A partial congruence
cover fixes $\N$ inside an arithmetic progression $\N=\N_0+kM$ and disposes of
most prime offsets for \emph{every} $k$; a residual set $U$ of offsets survives
and must be knocked out by luck, one candidate $k$ at a time. Under the
standard heuristic, the probability that a given candidate succeeds is
$e^{-E}$ with
\begin{equation}\label{eq:E}
E\;=\;\frac{|U|\cdot\mathrm{boost}}{\log\N},
\qquad
\mathrm{boost}\;=\;2\prod_{r\in\mathcal R}\frac{r}{r-1},
\end{equation}
so that $e^{E}$ candidates must be examined per expected hit. We call $E$ the
\emph{failure exponent}. It is the currency of the entire subject: every design
decision---which moduli to use, which residues, how many offsets to leave
residual, what digit budget to spend---is a trade in units of $E$, and every
piece of engineering (faster tests, better sieving, selective testing, more
hardware) buys a fixed number of nats of $E$ and no more.

Equation~\eqref{eq:E} is a heuristic, so we test it rather than assume it.
Across three independent covers and $1.2\times10^{9}$ scanned candidates, the
measured exponent $\Ehat$ agrees with the predicted $E$ to within $0.1$ nats
(Chapter~\ref{ch:engineering}). That agreement is what licenses the rest of the
thesis: once the model is trusted, questions about feasibility become
arithmetic.

\section{What is new here}

Chapters~\ref{ch:foundations} and~\ref{ch:engineering} are largely expository
and engineering: they set out the classical constructions, the covering
formulation, the calibrated cost model, and the search engine, and they report
the certified records and---equally important---the campaigns that failed.
The contributions begin in Chapter~\ref{ch:frontier}.

\begin{enumerate}
\item \textbf{The difficulty is an integrality gap} (\S\ref{sec:lp}). The
natural linear-programming relaxation of the covering problem is
\emph{vacuous}: fractionally, the available moduli cover $100.00\%$ of the
prime offsets at every target we tested. Consequently the whole failure
exponent is the price of \emph{rounding}: naive independent rounding leaves
$15.6\%$ of offsets uncovered, adaptive greedy leaves $7.56\%$, and the
fractional optimum leaves none. Every absolute percentage point recovered is
worth about $2.7$ nats, i.e.\ a factor of $15$ in compute.

\item \textbf{The rounding gap resists optimisation} (\S\ref{sec:door1}). We
attack it with exact large-neighbourhood search (certified optimal by LP
duality on $21$ of $22$ windows), with belief propagation and decimation, and
with a simultaneous exact solve of half the cover at once. The greedy solution
$|U|=867$ is never beaten. A random-restart probe shows why: the landscape is a
field of mutually distant, locally-exact-stable basins, and the adaptive greedy
construction lands about $100$ offsets deeper than a typical one.

\item \textbf{Sub-$100$-digit deserts reduce to a hard problem}
(Chapter~\ref{ch:barriers}). Oversized covers would make sub-$100$-digit
deserts easy if their Chinese-remainder representatives could be made small;
that requirement is exactly a modular subset-sum problem with per-position
choices at density $\approx1.026$, the regime where neither lattice reduction
nor birthday methods work. We also verify against the source that the relevant
CRT list-decoding radius is a factor $5.2$ short of what the construction would
need, and confirm experimentally that lattice methods do not penetrate the
Johnson-type radius even when solutions are abundant.

\item \textbf{A closure theorem} (Chapter~\ref{ch:closure}). Within a concrete
language of certificate schemas, no construction beats covering. Every family
closes for its own quantifiable reason---a $\sqrt{Q}$ value-set ceiling, a
purchase plan running $2.5\times$ over the design-entropy budget, a cluster
bound, lattice isolation, or an empty intersection at record scale---and the
theorem is conditional on two named hypotheses, one of which we pose as an open
problem.

\item \textbf{A methodological finding} (\S\ref{sec:referee}). We stress-tested
the closure theorem with an adversarial evolutionary search against a
machine-checkable scorer. Its first success was against the \emph{referee}, not
the theorem: it found a $50\times$ accounting error in our evaluator and
exploited it within ten candidates. The audit protocol caught it before any
claim was made, and repairing the referee strengthened the mathematics.
\end{enumerate}

\section{Notation and conventions}

Throughout, $\log_k$ is the $k$-fold iterated logarithm (so
$\log_2\N=\log\log\N$), and each extra $\log$ grows \emph{extremely} slowly;
$\theta(x)=\sum_{p\le x}\log p$ is Chebyshev's function;
$\gamma\approx0.5772$ is Euler's constant; $f\gg g$ means $f\ge cg$ for some
$c>0$ and all large arguments; $\pi(x)$ is the prime-counting function; and
``BPSW'' denotes the Baillie--PSW probable-prime test~\cite{PSW,BW}. All
entropies, costs and exponents are measured in \emph{nats} (natural logarithm
units), because the natural unit of a congruence modulo $r$ is $\log r$ and the
natural unit of a search is $\log(\text{candidates})$.

We are careful to separate three claims about any proposed record, and the
reader should hold us to the distinction throughout:
\begin{itemize}
\item[(a)] no prime below the stated value is a Goldbach summand of $\N$;
\item[(b)] $\N$ admits a Goldbach representation at all;
\item[(c)] the least summand \emph{equals} the stated value.
\end{itemize}
A congruence or a failed probable-prime test settles (a) unconditionally---a
prime can never fail such a test, so negative results are never in doubt. Claim
(b) needs an exhibited prime complement, and (c) needs both together with a
primality certificate for that complement. Where a statement is heuristic,
conditional, or merely measured, we say so: results proved at the level of
rigour of this thesis are marked as such, statements resting on standard but
unproved analytic heuristics are flagged, and purely empirical claims are
reported with the size of the experiment that supports them.

\chapter{Foundations: from factorials to covers, unconditionally}
\label{ch:foundations}

\begin{chapterabstract}
This chapter is the foundations paper, subsuming and updating the July
2026 manuscript \emph{Cute Goldbach Gaps}. After an elementary factorial
warm-up, we prove unconditionally that $\g$ is unbounded: a primorial
obstruction combined with Dirichlet's theorem produces, for every prime
$P$, an even $\N$ with an explicit partition and
$\g(\N)=\mathrm{nextprime}(P)$. We then formulate the construction as a
congruence-cover problem on the prime offsets, in complete and partial
forms, prove the realisation theorem that upgrades any cover to an
actual Goldbach number, and transfer the
Ford--Green--Konyagin--Maynard--Tao long-gap covering into an
unconditional lower bound on the maximal order of $\g$. The chapter
closes with the first-generation records this machinery produced --- a
$1020$-digit complete-cover baseline, a $237$-digit integer with
$\g=109\,621$, and a $199$-digit integer with $\g=105\,667$ --- which are
the incumbents every later chapter improves.
\end{chapterabstract}

\section{Factorial neighbourhoods and Wilson obstructions}\label{sec:factorial}

\begin{lemma}[Factorial composites]\label{lem:fact}
For $n\ge2$ and $2\le k\le n$, the number $n!+k$ is composite. For
$n\ge4$ and $2\le k\le n$, the number $n!-k$ is also composite.
\end{lemma}
\begin{proof}
In both cases $k\mid n!$, so $k\mid(n!\pm k)$. For the plus side
$n!+k>k$, so $n!+k$ is a proper multiple of $k$. For the minus side,
$n!-k=k\bigl(n!/k-1\bigr)$ with $n!/k-1\ge2$ exactly when $n!\ge3k$,
which holds for $n\ge4$; it fails at $n=3,\,k=3$, where $3!-3=3$ is
prime.
\end{proof}

\begin{lemma}[Wilson obstructions]\label{lem:wilson}
Let $n\ge2$. \emph{(A)} If $n+1$ is prime, then $(n+1)\mid n!+1$.
\emph{(B)} If $n+2$ is prime, then $(n+2)\mid n!-1$.
\end{lemma}
\noindent Recall \emph{Wilson's theorem}: a number $p>1$ is prime if and
only if $(p-1)!\equiv-1\pmod p$, i.e.\ the product $1\cdot2\cdots(p-1)$
is one less than a multiple of $p$. We use only the forward direction.
\begin{proof}
For (A) put $p=n+1$; Wilson's theorem gives $n!=(p-1)!\equiv-1\pmod p$.
For (B) put $p=n+2$; then $(p-1)!=(n+1)\,n!$ and $n+1\equiv-1\pmod p$,
so $-n!\equiv-1$, i.e.\ $n!\equiv1\pmod p$.
\end{proof}

\begin{corollary}\label{cor:facprime}
If $n!+1>n+1$ is prime, then $n+1$ is composite. In particular the
factorial primes $n!+1$ with $n\ge3$ (for example $n=3,11,27,37,41,\dots$)
occur only when $n+1$ is composite.
\end{corollary}

\subsection{The worked example \texorpdfstring{$\N=27!+28$}{N=27!+28}}

Take $\N=27!+28$, even with $29$ digits. Since $27!\equiv0\pmod r$ for
every prime $r\le27$,
\begin{equation}\label{eq:shift}
\N-q\;=\;27!+(28-q)\;\equiv\;28-q\pmod r\qquad(r\le27\ \text{prime}),
\end{equation}
so $\N-q$ is killed by a small prime $r$ exactly when $r\mid(28-q)$: a
shifted sieve that eliminates the class $q\equiv28\pmod r$.

\begin{proposition}\label{prop:71}
For $\N=27!+28$ every prime $q$ with $2\le q\le67$ has $\N-q$ composite,
while $\N-71$ is prime; hence $\g(\N)=71$.
\end{proposition}
\begin{proof}[Proof by explicit witnesses]
Table~\ref{tab:witness} gives, for each prime $q\le71$, a factor of
$\N-q$. For $q\le23$ the witness divides $28-q$ via~\eqref{eq:shift}.
Two cases below the winner are notable: $q=29$ is killed by $29$ because
$29=27+2$ is prime, so Lemma~\ref{lem:wilson}(B) gives
$29\mid27!-1=\N-29$; and $q=59$ survives the sieve (as $59-28=31$ is
$27$-rough) yet $41\mid27!-31=\N-59$. Finally $71-28=43$ is $27$-rough
and $27!-43$ is prime.
\end{proof}

\begin{table}[h]
\centering\small
\caption{Why every small prime fails as a summand of $\N=27!+28$.}\label{tab:witness}
\begin{tabular}{rlc@{\qquad}rlc}
\toprule
$q$ & $\N-q$ & factor & $q$ & $\N-q$ & factor\\
\midrule
$2$&$27!+26$&$2$ & $31$&$27!-3$&$3$\\
$3$&$27!+25$&$5$ & $37$&$27!-9$&$3$\\
$5$&$27!+23$&$23$ & $41$&$27!-13$&$13$\\
$7$&$27!+21$&$3$ & $43$&$27!-15$&$3$\\
$11$&$27!+17$&$17$ & $47$&$27!-19$&$19$\\
$13$&$27!+15$&$3$ & $53$&$27!-25$&$5$\\
$17$&$27!+11$&$11$ & $59$&$27!-31$&$\mathbf{41}$ (generic)\\
$19$&$27!+9$&$3$ & $61$&$27!-33$&$3$\\
$23$&$27!+5$&$5$ & $67$&$27!-39$&$3$\\
$29$&$27!-1$&$\mathbf{29}$ (Wilson) & $71$&$27!-43$&\textbf{prime}\\
\bottomrule
\end{tabular}
\end{table}

By~\eqref{eq:shift} a prime $q$ can even be \emph{tested} only if $q-28$
is $27$-rough, so the viable candidates $29,59,71,89,\dots$ are sparse,
and $71$ is merely the first with a prime complement. Sharper still:
$27!+1$ \emph{is} prime (Corollary~\ref{cor:facprime}), the natural
bridge $q=27$, but $27$ is composite, so one falls through to $71$. For
indices that are simultaneously prime and factorial-prime indices
($n=3,11,37,41$) the bridge works and $\g(n!+(n+1))=n$. Heuristically,
the first viable offsets in this factorial family lie on the scale of
$n$, corresponding to roughly $\log\N/\log\log\N$; the rest of the
chapter produces much larger values.

\section{Primorial deserts, made unconditional}\label{sec:primorial}

The factorial wastes each small prime on a residue class. Taking $\N$ to
be a multiple of a primorial wastes nothing.

\begin{lemma}[Primorial core]\label{lem:primcore}
Let $P$ be prime, let $m_P=\prod_{p\le P}p$ be the \emph{primorial} ---
the product of every prime up to $P$ (for instance
$m_{11}=2\cdot3\cdot5\cdot7\cdot11=2310$) --- and let $\N$ be a multiple
of $m_P$ with $\N>2P$. Then $\N-q$ is composite for every prime
$q\le P$; i.e.\ no prime $\le P$ is a Goldbach summand of $\N$.
\end{lemma}
\begin{proof}
For a prime $q\le P$ we have $q\mid m_P\mid\N$ (as $q$ is one of the
factors of $m_P$), hence $q\mid\N-q$, and $\N-q>P\ge q$, so $\N-q$ is a
proper multiple of $q$ and therefore composite.
\end{proof}

Lemma~\ref{lem:primcore} alone gives only claim~(a) of the standard of
evidence. To reach unboundedness with an \emph{actual} Goldbach number
we add Dirichlet's theorem.

\begin{theorem}[Unconditional unboundedness]\label{thm:dirichlet}
For every prime $P$ put $s=\mathrm{nextprime}(P)$. There are infinitely
many primes $t\equiv-s\pmod{m_P}$, and for any such $t$ the even integer
$\N=s+t$ satisfies $m_P\mid\N$, has the explicit Goldbach partition
$\N=s+t$, and has
\[
\g(\N)\;=\;\mathrm{nextprime}(P)\;>\;P.
\]
Hence $\sup_\N\g(\N)=\infty$ unconditionally. Moreover the least such
$t$ satisfies $\log\N\ll P$ by Linnik's theorem~\cite{Linnik}, so
$\g(\N)\gg\log \N$ along this family.
\end{theorem}
\begin{proof}
As $\gcd(s,m_P)=1$ (because $s>P$), Dirichlet~\cite{Dirichlet} gives
infinitely many primes $t\equiv-s\pmod{m_P}$. Then
$\N=s+t\equiv0\pmod{m_P}$ and $\N$ is even. By Lemma~\ref{lem:primcore}
no prime $\le P$ is a summand; the only primes below
$s=\mathrm{nextprime}(P)$ are those $\le P$, so the least summand is at
least $s$. But $\N-s=t$ is prime, so $s$ is a summand; hence
$\g(\N)=s$. Linnik's theorem bounds the least such prime by
$t\ll m_P^{L}$ for an absolute constant $L$, and
$\log m_P=\theta(P)\sim P$, so $\log\N\ll P\sim s=\g(\N)$.
\end{proof}

\begin{remark}[Concrete instances]\label{rem:concrete}
For $P=13$ one has $m_{13}=30030$, $s=17$, and the least prime
$t\equiv-17\pmod{30030}$ is $t=30013$, giving
$\N=30030=2\cdot3\cdot5\cdot7\cdot11\cdot13$ with partition
$30030=17+30013$ and $\g(\N)=17$ --- all primality deterministic. A
larger fully proven example is the primorial
$m_{43}=13082761331670030$ itself: its minimal partition is
$m_{43}=89+13082761331669941$ (both proven prime, being $<2^{64}$), so
$\g(m_{43})=89$ unconditionally.
\end{remark}

\begin{remark}[The roughness bonus]\label{rem:rough}
For the bare choice $\N=m_P$, a prime $q>P$ makes $\N-q$ coprime to all
primes $\le P$ --- the complement is \emph{$P$-rough} --- which improves
its chances of being prime, since the small primes that catch most
composites are already excluded. By \emph{Mertens' theorem} (which says
$\prod_{p\le P}(1-1/p)\sim e^{-\gamma}/\log P$: the product of $1-1/p$
over primes $p\le P$ shrinks like $1/\log P$), this sieving effect
multiplies the heuristic primality probability by
$\prod_{p\le P}(1-1/p)^{-1}\sim e^\gamma\log P$. The expected least
surviving summand then lands near $P(1+e^{-\gamma})\approx1.56\,P$.
Either way the primorial pays a modulus $m_P$ with
$\log m_P=\theta(P)\sim P$ --- exponentially large in the target --- and
the next section removes this waste. The roughness bonus itself,
however, is not waste but the engine of everything that follows: it
reappears in Chapter~\ref{ch:costmodel} as the ``boost'' factor of the
cost model.
\end{remark}

\section{Prime-offset congruence covers}\label{sec:covering}

\subsection{The construction}

A primorial spends one prime per summand; a prime-offset congruence
cover reuses each prime on a whole residue class. The one structural
constraint is parity: a sum of two odd primes is even, so $\N$ is kept
even and the prime $2$ kills only $q=2$; every \emph{odd} prime below
the target must be covered by odd sieving primes.

The strategy itself is classical. Covering systems of congruences were
introduced by Erd\H{o}s precisely to obstruct an additive
representation: his construction of an infinite arithmetic progression
of odd integers not of the form $2^k+p$~\cite{Erdos1950} is the direct
ancestor of this section (and of the Sierpi\'nski--Riesel tradition that
grew from it). What changes here is only the target: the offsets to be
blocked are the primes below $Q$ rather than the powers of two.

\begin{proposition}[Covering certificate]\label{prop:cover}
Let $Q\ge3$ and let $\{(r,a_r):r\in\mathcal R\}$ be a system of odd
primes $\mathcal R$ such that every odd prime $q<Q$ has
$q\equiv a_r\pmod r$ for some $r\in\mathcal R$. If $\N$ is even,
$\N\equiv a_r\pmod r$ for all $r\in\mathcal R$, and
\[
\N\;>\;Q+\max_{r\in\mathcal R}r,
\]
then $\N-q$ is composite for every prime $q<Q$; i.e.\ no prime $<Q$ is a
Goldbach summand of $\N$.
\end{proposition}
\begin{proof}
For $q=2$, $\N-2$ is even and $>2$. For an odd prime $q<Q$, pick
$r\in\mathcal R$ with $q\equiv a_r\equiv\N\pmod r$; then $r\mid\N-q$,
and $\N-q>(Q+\max_r r)-Q\ge r$, so $r\mid\N-q$ properly and $\N-q$ is
composite.
\end{proof}

\begin{remark}
The size hypothesis only has to make every divisor $r\mid\N-q$ proper,
for which $\N>Q+\max_r r$ suffices; the much stronger $\N>Q\prod_r r$ is
unnecessary. A common solution $\N$ exists modulo $M=2\prod_{r}r$ by the
Chinese Remainder Theorem --- the classical fact that congruences to
pairwise coprime moduli can be satisfied simultaneously, and that the
solutions form exactly one residue class modulo the product --- and any
large enough even representative qualifies.
\end{remark}

We build coverings greedily and deterministically.

\medskip
\noindent\fbox{\parbox{0.97\textwidth}{%
\textbf{Algorithm (greedy covering of the odd primes below $Q$).}
\begin{enumerate}\setlength\itemsep{1pt}
\item Let $T=\{\text{odd primes }<Q\}$ and require $\N\equiv0\pmod2$.
\item For $r=3,5,7,11,\dots$ in increasing order, while
$T\ne\varnothing$: let $a_r$ be a residue mod $r$ containing the most
elements of $T$ (smallest such residue on ties), require
$\N\equiv a_r\pmod r$, and delete from $T$ every $q\equiv a_r\pmod r$.
\item Recover $\N$ from the congruences by the Chinese Remainder
Theorem.
\end{enumerate}}}
\medskip

The resulting modulus is $M=2\prod_{r\in\mathcal R}r$ --- twice the
product of the sieving primes actually used. For the computed targets it
is far smaller than the primorial $\prod_{p<Q}p$, because each sieving
prime is reused across an entire residue class rather than assigned to
one offset. (Chapter~\ref{ch:records} replaces this naive greedy pass
with weighted set cover plus annealing; Chapter~\ref{ch:closure} asks,
and largely answers, how far from optimal such covers are.)

\subsection{Partial covers and progression search}\label{sec:partial}

A complete cover is convenient to verify, but it is not the cheapest way
to make $\N$ small. Near the end of the greedy algorithm, a new sieving
prime may remove only one or two remaining offsets while still
multiplying the modulus by $r$. That is an expensive way to buy a single
compositeness proof.

The alternative is to stop early. Let $U$ be the odd prime offsets below
$Q$ not covered by the chosen congruences --- the \emph{residual set}.
CRT gives one residue $\N_0$ modulo
\[
M=2\prod_{r\in\mathcal R}r,
\]
so every candidate has the form
\begin{equation}\label{eq:progression}
\N=\N_0+kM,\qquad k=0,1,2,\dots
\end{equation}
The fixed congruences dispose of most offsets for every $k$. For each
remaining $q\in U$, we test $\N-q$ directly. A failed probable-prime
test is a rigorous compositeness result: pseudoprimes can create false
positives, but they cannot turn a prime into a reported composite. Thus
a fast one-sided test is safe as a search filter, while the final
surviving complement still needs a primality certificate.

\begin{proposition}[Partial-cover certificate]\label{prop:partial}
Let $U$ be the odd primes $q<Q$ not covered by the congruences of
Proposition~\ref{prop:cover}. Let $\N$ be a sufficiently large even
solution of those congruences. If $\N-q$ is composite for every
$q\in U$, then no prime $q<Q$ is a Goldbach summand of $\N$.
\end{proposition}
\begin{proof}
The congruence witnesses prove compositeness for every covered odd
prime; parity handles $q=2$; and the hypothesis handles the residual
set $U$.
\end{proof}

Most values of $k$ fail the hypothesis --- some residual complement
turns out to be prime, which is not a disaster but a small Goldbach
partition, i.e.\ a losing lottery ticket --- and the construction
becomes a \emph{search} along the progression \eqref{eq:progression}.
What that search costs, and how predictably, is the subject of
Chapter~\ref{ch:costmodel}; it is the single question the rest of the
thesis orbits.

\subsection{Realising a covering unconditionally}

Proposition~\ref{prop:cover} still only gives claim~(a). The same
Dirichlet idea as in \S\ref{sec:primorial}, applied twice, upgrades any
covering to an actual Goldbach number.

\begin{theorem}[Unconditional realisation]\label{thm:coverdirichlet}
Let $Q\ge3$ and let $\{(r,a_r):r\in\mathcal R\}$ cover the odd primes
below $Q$, where every $r\in\mathcal R$ is an odd prime. Then there are
infinitely many even integers $\N>Q+\max_{r\in\mathcal R}r$ admitting an
explicit Goldbach partition $\N=s+t$ with $\N\equiv a_r\pmod r$ for all
$r\in\mathcal R$. Every such $\N$ has $\g(\N)\ge Q$.
\end{theorem}
\begin{proof}
For each $r\in\mathcal R$ choose
$b_r\in\{1,\dots,r-1\}\setminus\{a_r\}$, and also require
$s\equiv1\pmod2$. The CRT gives a reduced residue class $b$ modulo
$2\prod_r r$. Dirichlet gives infinitely many primes in that class; call
one of them $s$. If $s<Q$, then $s$ is an odd covered prime, so
$s\equiv a_r\pmod r$ for some $r$, contradicting
$s\equiv b_r\not\equiv a_r\pmod r$. Hence $s\ge Q$.

Next require
\[
t\equiv a_r-s\pmod r\quad(r\in\mathcal R),\qquad t\equiv1\pmod2.
\]
Every residue $a_r-s$ is nonzero modulo $r$, so this is a reduced
residue class modulo $2\prod_r r$. Dirichlet gives infinitely many prime
values of $t$; keep those large enough that $s+t>Q+\max_r r$, and put
$\N=s+t$. Then $\N$ is even, $\N\equiv a_r\pmod r$, and $\N=s+t$ is an
explicit Goldbach partition. Proposition~\ref{prop:cover} excludes every
prime below $Q$, so $\g(\N)\ge Q$.
\end{proof}

\begin{remark}[Quantitative form for initial-segment covers]\label{rem:size}
The covering used in Theorem~\ref{thm:lower} below has one modulus for
every prime $p\le x$. In that setting the two residue classes above are
reduced modulo $M=2\prod_{p\le x}p$. Linnik's theorem~\cite{Linnik}
supplies primes $s,t\ll M^L$ for an absolute constant $L$. Their
residues are nonzero modulo every prime at most $x$, and the covering
property forces the odd prime $s$ to satisfy $s\ge Q$. Consequently the
resulting integer satisfies
\[
\log\N\ll\log M=\theta(x)\ll x.
\]
This quantitative statement is used only for those initial-segment
covers; no such bound is asserted for an arbitrary sparse finite cover.
\end{remark}

\section{Transfer from large prime gaps}\label{sec:gaptransfer}

\begin{proposition}\label{prop:transfer}
If $[\N-Q,\N-2]$ contains no prime, then no prime $q\le Q$ is a Goldbach
summand of $\N$. Consequently, if $\N$ admits a Goldbach representation,
then $\g(\N)>Q$.
\end{proposition}
\begin{proof}
Every prime $q\in[2,Q]$ has $\N-q\in[\N-Q,\N-2]$, hence $\N-q$ is
composite.
\end{proof}

So a prime gap of length $\ge Q$ just below $\N$ forces $\g(\N)>Q$:
\emph{prime gaps transfer to least Goldbach summands}. The transfer is,
however, strictly one-directional. Forcing $\g(\N)>Q$ requires $\N-q$
composite only for \emph{prime} offsets $q<Q$, not for every integer in
an interval; the covering of Proposition~\ref{prop:cover} covers just
the primes below $Q$, a far sparser target. Thus ``large least Goldbach
summand'' is a relaxation of ``large prime gap,'' and one should not
expect the two problems to coincide. Chapter~\ref{ch:gaps} measures the
separation on certified examples; here we exploit the easy direction.

Writing $G(X)$ for the largest gap between consecutive primes below
$X$, the work of Ford--Green--Konyagin--Tao and of
Maynard~\cite{FGKT,Maynard}, sharpened in~\cite{FGKMT}, gives
\[
G(X)\;\gg\;\frac{\log X\,\log_2 X\,\log_4 X}{\log_3 X},
\]
improving Rankin~\cite{Rankin}. In words: there are stretches of
consecutive composite integers below $X$ of length at least a constant
times $\log X\cdot\log_2 X\cdot\log_4 X/\log_3 X$. The dominant factor
is $\log X$ (the average prime gap near $X$); the three iterated-log
factors are a slowly growing extra boost, and since each $\log_k$ grows
glacially this is only a little more than $\log X$. Feeding the covering
that underlies these bounds into Theorem~\ref{thm:coverdirichlet} (which
costs only a constant factor in $\log\N$, by Linnik) makes the transfer
unconditional:

\begin{theorem}\label{thm:lower}
There are infinitely many even integers $\N$, each with an explicit
Goldbach representation, such that
\[
\g(\N)\;\gg\;\frac{\log \N\,\log_2 \N\,\log_4 \N}{\log_3 \N}.
\]
\end{theorem}

\begin{proof}
The Erd\H{o}s--Rankin method, in the sharpened form of
Ford--Green--Konyagin--Maynard--Tao, supplies the covering. Precisely,
\cite{FGKMT} shows (in the standard covering formulation of the method,
the quantity $Y(x)$ of \cite{FGKT,FGKMT}) that there is an absolute
constant $c>0$ such that for all large $x$ one may choose residue
classes $a_p\pmod p$, one for each prime $p\le x$, whose union contains
every integer in $[1,y]$, where
\[
y=\Bigl\lfloor c\,\frac{x\,\log x\,\log_3 x}{\log_2 x}\Bigr\rfloor;
\]
the prime-gap bound is deduced from exactly this covering by the Chinese
Remainder Theorem.

\emph{Normalisation.} If $a_2$ is odd, replace every class $a_p$ by
$a_p-1$; the shifted system covers $[0,y-1]\supseteq[1,y-1]$. So, at the
cost of replacing $y$ by $y-1$, we may assume $2\mid a_2$. Then no odd
integer in $[1,y-1]$ is covered by the class modulo $2$, so every odd
integer $m\le y-1$ --- in particular every odd prime $q<Q:=y-1$ ---
satisfies $q\equiv a_r\pmod r$ for some \emph{odd} prime $r\le x$. The
pairs $(r,a_r)$ over the odd primes $r\le x$ therefore form a covering
of the odd primes below $Q$ in the sense of
Proposition~\ref{prop:cover}: it covers far more (all odd integers below
$Q$), but only the primes are needed, and self-covering ($r=q$,
$a_r=0$) is harmless, exactly as in the computed systems of
Chapter~\ref{ch:records}.

\emph{Realisation.} Theorem~\ref{thm:coverdirichlet} applied to this
covering yields infinitely many even $\N$ with an explicit partition and
$\g(\N)\ge Q$; by Remark~\ref{rem:size} we may take
\[
\log\N\;\ll\;\log2+\theta(x)\;\ll\;x.
\]

\emph{Conclusion.} Let $C$ be the absolute constant with
$\log\N\le Cx$. For large $u$ the function
$u\mapsto u\log u\,\log_3u/\log_2u$ is increasing, and each iterated
logarithm of $\log\N/C$ is $(1+o(1))$ times the corresponding iterated
logarithm of $\log\N$; hence
\begin{align*}
\g(\N)\;\ge\;Q\;&\gg\;\frac{x\,\log x\,\log_3x}{\log_2x}
\;\ge\;\frac{\tfrac{\log\N}{C}\,\log\bigl(\tfrac{\log\N}{C}\bigr)\,\log_3\bigl(\tfrac{\log\N}{C}\bigr)}{\log_2\bigl(\tfrac{\log\N}{C}\bigr)}\\[2pt]
&\gg\;\frac{\log\N\,\log_2\N\,\log_4\N}{\log_3\N}.
\end{align*}
For each sufficiently large $x$, select one realisation $\N_x$
satisfying this Linnik bound. Since $\g(\N_x)\ge Q(x)\to\infty$, the
selected integers are unbounded and therefore contain infinitely many
distinct values. This proves the claim.
\end{proof}

Together with the Granville--van de Lune--te Riele
conjecture~\eqref{eq:GLR}, this brackets the maximal order of $\g$
between a logarithmic-type lower bound and a $(\log\N)^2$-type
conjectural ceiling:
\[
\underbrace{\frac{\log \N\,\log_2 \N\,\log_4 \N}{\log_3 \N}}_{\text{rigorous lower-bound scale}}
\;\ll\;\max_{\N'\le \N}\g(\N')
\;\lesssim\;
\underbrace{(\log \N)^2\log\log \N}_{\text{conjectural ceiling}}.
\]
The explicit constructions of this thesis do not settle that gap; their
value is that they turn the problem into auditable objects --- a
congruence list, a CRT progression, a finite residual set, and a
primality certificate --- and, in Chapters~\ref{ch:limits}
and~\ref{ch:closure}, into measured statements about which parts of the
gap are closable at all. We note one reappearance of this asymptotic
machinery: the FGKMT covering enters again in
Chapter~\ref{ch:closure} not as a lower bound but as the one known
\emph{rounding technology} that might beat greedy covers, which is this
thesis's central open problem.

\section{The first-generation records}\label{sec:firstgen}

The machinery above produced the first constructive records in this
problem family, which the rest of the thesis treats as its baseline.
All three claims (a)/(b)/(c) were discharged to the standard of
Chapter~\ref{ch:intro}, and all were re-verified inside the thesis
pipeline before any successor record was accepted.

\paragraph{The complete-cover baseline.} A complete greedy cover of the
odd primes below $Q=10^5$ uses $356$ congruences and forces a modulus of
$1020$ digits; its least even representative has $\g(\N)=100\,747$.
Every excluded offset has a small factor visible directly from the
congruence system, so verification is little more than modular
arithmetic --- a \emph{proof-oriented build} that overpays roughly
sevenfold in digits for certainty it does not need. (A complete cover at
$Q=10^8$ was also computed as a scaling probe: $327\,455$ digits, with
no partition exhibited.)

\paragraph{The hybrid ablation.} Stopping the cover early and searching
the progression \eqref{eq:progression} shrinks the integer dramatically.
The engineering path, at fixed target $Q=10^5$:

\begin{table}[h]
\centering\small
\caption{First-generation ablation at target $Q=10^5$. ``Congr.''
includes parity; ``residual'' counts prime offsets below $10^5$ left
for direct testing; the last column is the first prime offset whose
complement is (certified) prime.}\label{tab:ablation}
\begin{tabular}{lrrrrrr}
\toprule
Method & Congr. & Digits $M$ & Residual & $k$ & Digits $\N$ & $\g(\N)$\\
\midrule
Complete greedy cover & 356 & 1020 & 0 & 0 & 1020 & $100\,747$\\
Prefix partial cover & 168 & 416 & 385 & 74 & 418 & $101\,599$\\
Coordinate descent & 110 & 248 & 580 & $1\,751$ & 251 & $100\,981$\\
Sparse weighted cover & 103 & 232 & 610 & $299\,948$ & 237 & $109\,621$\\
\bottomrule
\end{tabular}
\end{table}

\paragraph{The $237$-digit height record.} The final row is the
first-generation height incumbent: a $237$-digit even integer
$\N_{237}=\N_0+299\,948\,M$ over a partial cover of $102$ odd
congruences plus parity ($M$ of $232$ digits, largest modulus $1217$),
with
\[
\g(\N_{237})=109\,621.
\]
Every prime offset below $109\,621$ is excluded --- $9\,691$ by explicit
congruence divisors and $731$ by strong base-2 compositeness witnesses
--- and the complement $\N_{237}-109\,621$ is ECPP-certified prime.

\paragraph{The $199$-digit dual incumbent.} Re-optimising the cover for
the strict budget $\N<10^{199}$ gives a $199$-digit integer
$\N_{199}=\N_0+876\,554\,M$ ($89$ odd congruences plus parity, $M$ of
$193$ digits) with
\[
\g(\N_{199})=105\,667,
\]
proved by $9\,335$ parity-or-congruence witnesses, $744$ strong base-2
witnesses, and a $34$-step elliptic-curve certificate for the
complement. It entered the thesis as the incumbent for \emph{both}
bounded games: $T(100\,000)\le\N_{199}$ and $R(10^{199})\ge105\,667$.

\medskip
These records made the covering paradigm concrete, and they framed the
questions the thesis then had to answer. Why $237$ digits and not $150$?
Is the greedy cover any good? What does a progression search actually
cost, and can that cost be predicted rather than experienced? The next
chapter builds the measuring instruments; Chapter~\ref{ch:records}
spends them.

\chapter{The cost of a desert: a calibrated model and an engine}\label{ch:engineering}

Chapter~\ref{ch:foundations} reduced desert construction to a search: fix a
partial cover, walk the progression $\N=\N_0+kM$, and test the residual offsets.
This chapter answers the obvious next question---how long does that take?---and
then reports what happened when we ran it. The answer is a single exponent, and
the exponent is accurate enough that most of the rest of this thesis can proceed
by arithmetic rather than by experiment.

\section{The failure exponent}

Fix a partial cover with modulus $M=2\prod_{r\in\mathcal R}r$ and residual set
$U$ of prime offsets below $Q$. For a candidate $\N=\N_0+kM$, the covered
offsets are composite by construction; the candidate succeeds precisely when
\emph{every} residual complement $\N-q$, $q\in U$, is also composite.

Two effects set the probability. First, each complement is a number of size
about $\N$, so by the prime number theorem it is prime with probability about
$1/\log\N$. Second---and this is the roughness bonus of
Remark~\ref{rem:rough}---the cover has already forced $\N\equiv a_r$, so each
complement $\N-q$ avoids divisibility by every $r\in\mathcal R$ (for the
offsets that were \emph{not} killed by $r$), which raises its chance of being
prime. Mertens' theorem turns that into a multiplicative correction. Writing
\begin{equation}\label{eq:boost}
\mathrm{boost}\;=\;2\prod_{r\in\mathcal R}\frac{r}{r-1},
\end{equation}
the probability that one residual complement is prime is
$\approx\mathrm{boost}/\log\N$, and, treating the $|U|$ residual complements as
independent, the probability that a candidate survives all of them is
\[
\Bigl(1-\frac{\mathrm{boost}}{\log\N}\Bigr)^{|U|}\;\approx\;e^{-E},
\qquad
E\;=\;\frac{|U|\cdot \mathrm{boost}}{\log\N}.
\]
Thus $e^{E}$ candidates must be scanned per expected success. The independence
step is a heuristic---it is the Hardy--Littlewood picture applied to a fixed
arithmetic progression---and Chapter~\ref{ch:closure} shows that it is in a
precise sense the \emph{only} heuristic in play, because any correlation between
two complements is forced to act through small shared divisors, which the cover
already controls.

The formula explains the shape of the whole subject. Three levers move $E$:
\begin{itemize}
\item \textbf{More moduli} shrink $|U|$ but inflate $\mathrm{boost}$ and
$\log M$. The balance between these is not obvious, and \S\ref{sec:equilibrium}
shows it is an equilibrium rather than a budget constraint.
\item \textbf{A bigger $\N$} increases $\log\N$ and so lowers $E$ directly. This
is why the height game is easy and the budget game is hard: digits are the
cheapest resource in the problem.
\item \textbf{A higher target $Q$} raises $|U|$ roughly in proportion to
$\pi(Q)$, which is why record values of $\g$ get expensive fast.
\end{itemize}

\section{Calibration: does the model hold?}\label{sec:calibration}

A cost model that has never been checked is a wish. We therefore instrumented
every search to report, per slice of the progression, the measured pass rate
$\hat p$ of the probable-prime tests, the mean number of surviving offsets per
candidate, and the implied exponent
\[
\Ehat\;=\;\frac{(\text{mean survivors per candidate})\times\hat p}{1},
\]
which is the empirical analogue of $E$.

Across three independent covers and roughly $1.2\times10^{9}$ scanned
candidates, we find
\[
|\Ehat-E|\;\le\;0.1\ \text{nats},
\]
with per-slice values stable to within a couple of percent over hundreds of
slices. Two consequences matter. The obvious one is that projections in later
chapters can be made by arithmetic. The subtler one is a genuine bound on
hidden structure: if some unexploited arithmetic correlation made constructed
$\N$ systematically luckier than the model predicts, it would show up as a
persistent gap between $\Ehat$ and $E$. At $1.2\times10^{9}$ candidates, any
such effect is worth at most $0.1$ nats, a factor of $1.11$. We use this as a
quantitative constraint in Chapter~\ref{ch:closure}, where it rules out an
entire class of would-be mechanisms.

\section{The engine}\label{sec:engine}

The search is embarrassingly parallel and utterly dominated by one primitive: a
base-$2$ Fermat test on a number of the size of $\N$. Everything else is
bookkeeping designed to reduce the number of such tests.

\paragraph{Sieving.} For each residual offset $q$ and each small prime
$s\nmid M$, the candidates $k$ for which $s\mid\N-q$ form an arithmetic
progression in $k$, computable from one modular inverse. Marking those
progressions removes most $(k,q)$ pairs before any expensive test. The optimal
sieve bound is empirical: at $199$ digits we measure $B=10^{6}$ to beat
$B=3\cdot10^{6}$ ($901$ versus $816$ thousand candidates per second), because
beyond that the sieve costs more than the tests it saves.

\paragraph{Amortisation across variants.} A cover admits many
\emph{residue-swap variants}: reassigning some $a_r$ produces a different
progression with the same modulus $M$ and the same $|U|$, hence the same $E$.
These are free additional lottery tickets. Because they share $M$, they also
share every modular inverse in the sieve, so caching that table across variants
gave a measured $+47\%$ in throughput.

\paragraph{Survivor-ordered testing.} Within a block, candidates differ in how
many residual offsets survived the sieve. A candidate with $a$ survivors
succeeds with probability $(1-p)^{a}$, so hits concentrate sharply in the
low-survivor tail. Testing candidates in ascending order of $a$ is
result-preserving and measurably accelerates \emph{discovery} (a factor of
$1.83$ in expected time-to-first-hit at record scale).

\paragraph{Selective testing.} The natural next step is to \emph{skip} the
crowded candidates entirely. We implemented this with exact accounting: the
engine computes the survivor histogram, weights each candidate by
$(1-\hat p)^{a}$, and reports the fraction of total hit mass retained. This
keeps the model honest---$\hat p$ remains unbiased, so calibration still
works---while allowing a principled speed/coverage trade. At $199$ digits,
skipping the fullest $92\%$ of each block took the raw scan rate from
$1\,772$ to about $16\,300$ thousand candidates per second while retaining
$27.6\%$ of the hit mass: a net discovery rate about $3\times$ better.

The reason this works is worth stating, because it is not obvious. Each
candidate dies at its \emph{first} passing complement, so the expected number of
Fermat tests per candidate is about $1/\hat p\approx17$ regardless of how many
survivors it has. Testing cost is therefore nearly flat in $a$, while success
probability decays geometrically in $a$---so discarding high-$a$ candidates is
almost pure profit.

\paragraph{Hardware.} The same algorithm on GPU is a different animal. A CUDA
implementation using a bit-matrix sieve and Montgomery CIOS Fermat waves,
validated bit-for-bit against the reference implementation, sustains about
$2.5\times10^{6}$ tests/s on a T4 and $8.5\times10^{6}$ on an A100, versus
roughly $6\times10^{3}$ on four CPU cores. That ratio, about $430$, is
$\log 430\approx6.1$ nats---and the E-gap between the rungs the two platforms
can reach is likewise about $6$. The record frontier is logarithmic in compute,
and this is the cleanest illustration of it in our data: covers and strategy
were already saturated at CPU scale, and the GPU rung was bought purely with
throughput.

\section{Search strategy: depth, breadth, and the ratchet}\label{sec:ratchet}

Given a fixed compute budget, how should it be spent? We ran the experiment
both ways and the answer was unambiguous.

\emph{Breadth-first} designs scan many residue-swap variants shallowly. Because
the attainable $\N$ stays close to $\N_0$, every hit is small---good for the
threshold game---but the design has a hard ceiling: when the variant pool is
exhausted, the campaign ends with nothing.

\emph{Depth-first} designs scan one progression far. The runway is unbounded, and
because $\N$ grows only logarithmically with $k$, bad luck costs digits rather
than terminating the campaign. Our $150$-digit threshold record was found at
$k=6.21\times10^{10}$ of a full $[0,7.5\times10^{10})$ scan---at $63.8\%$
depth, a central draw against a cumulative expectation of $0.40$.

The synthesis is what we call the \textbf{ratchet}: design a cover, run it until
a hit, \emph{bank a certified record}, then re-tune the cover for the next rung
and repeat. At $Q\approx10^{5}$ this produced the descending ladder
$195\to179\to150$ digits, each rung a certified record in its own right. Under a
fixed budget the ratchet dominates a single grand-slam campaign for a structural
reason: it converts a lottery with one payout into a sequence of lotteries each
of which pays, and each rung also re-calibrates $\Ehat$ before more compute is
committed to the next.

\section{Results, including the ones that failed}\label{sec:campaigns}

Table~\ref{tab:records} listed the successes. Honesty requires the rest of the
distribution, so Table~\ref{tab:campaigns} reports the campaigns that produced
no record. In each case we give the survival probability---the model's own
prediction of how likely it was to see nothing, given the compute actually
spent---so the reader can judge whether the failure indicts the model or merely
the dice.

\begin{table}[h]
\centering\small
\caption{Campaigns that produced no record. ``Survival'' is $e^{-\lambda}$ where
$\lambda$ is the cumulative expected number of hits over the compute actually
spent; a value near $1$ means the campaign never had a real chance, a small
value means it was unlucky.}\label{tab:campaigns}
\begin{tabular}{lrrrl}
\toprule
Campaign & Target & Candidates & Survival & Outcome\\
\midrule
E10 (budget game, $Q=107\,720$) & $\g>107\,719$ & $2.4\cdot10^{8}$ & $0.17$ & unlucky, in-model\\
Campaign 2 ($Q=99\,991$, 168d cover) & $T(10^5)$ ladder & $5.6\cdot10^{8}$ & $0.22$ & unlucky, in-model\\
Campaign 3 ($Q=122\,000$) & $\g>119\,419$ & $4.0\cdot10^{8}$ & $0.93$ & stopped early\\
\bottomrule
\end{tabular}
\end{table}

Two of these were genuinely unlucky: a $0.17$ survival means the campaign had a
$83\%$ chance of producing a record and did not. That is a one-in-six outcome,
not a broken model, and in both cases the per-slice statistics tracked the
prediction throughout---which is precisely the discipline that lets us say so
with a straight face rather than reaching for an explanation. The third was
stopped by hand after $2.3\%$ of its budget when we learned that a parallel
effort was attacking the same target; its checkpoint is committed and it remains
resumable.

The methodological point is that a calibrated model turns a negative result into
information. Without $E$, three failed campaigns would be an embarrassment or an
excuse. With $E$, they are three consistent measurements of the same exponent,
and they are part of the $1.2\times10^{9}$-candidate calibration that
\S\ref{sec:calibration} rests on.

\section{Certification}\label{sec:certification}

Every record in this thesis is packaged so that its verification is independent
of its discovery. The pipeline is deliberately dull:

\begin{enumerate}
\item \textbf{Reconstruct.} From the committed cover and $(\N_0,M,k)$,
recompute $\N$ and check it against the stored decimal string.
\item \textbf{Exclude.} For every prime $q<\g$, produce a witness: either the
modulus $r$ with $r\mid\N-q$ (checked by division) or a strong base-$2$
compositeness witness (checked by modular exponentiation). Both are
unconditional; a prime cannot fail either test.
\item \textbf{Represent.} Verify that $\g$ is prime and that $\N-\g$ is prime,
the latter by an Atkin--Morain ECPP certificate~\cite{GoldwasserKilian,AtkinMorain}
generated by PARI/GP~\cite{PARI}.
\item \textbf{Re-verify independently.} A second checker, written against the
record file rather than the search code, re-derives every witness and validates
the certificate chain step by step---including the projective
elliptic-curve arithmetic---without calling the library that produced it.
\item \textbf{Hash.} A SHA-256 manifest fixes the artifacts.
\end{enumerate}

Step 4 is not ceremony. Re-verifying the two incumbent records at the start of
this programme is what established the baseline the rest of the work is measured
against, and the independent checker has since caught discrepancies in
intermediate artifacts that the primary pipeline reported as clean. Verification
is also far cheaper than discovery: it depends only on the final cover, the
finite residual check, and the certificate---never on reproducing the heuristic
path that found them.

\chapter{The frontier: where covering stops, and why}\label{ch:frontier}

With a calibrated exponent in hand we can ask sharp questions. How small can a
desert be made at a given target $Q$? How large can $\g$ be made under a given
digit budget? And---the question this chapter is really about---is the answer
determined by mathematics we could improve, or by a wall?

The answer turns out to be surprisingly clean, and surprisingly unlike what we
expected when we started. The difficulty of the whole subject is an
\emph{integrality gap}: fractionally, the covering problem is trivial. All of
the exponent, every last nat of it, is the price of rounding a fractional cover
to an integral one.

\section{The frontier, measured}

Table~\ref{tab:frontier} gives the achievable exponent as a function of the
digit budget, at fixed target $Q=10^{5}$: for each cap on the number of digits
of $\N$, we build the best cover we can and report $|U|$ and $E$.

\begin{table}[h]
\centering\small
\caption{Cover frontier at $Q\approx10^{5}$: the exponent achievable under a
digit cap. Fewer digits force a smaller modulus, hence fewer moduli, hence a
larger residual set \emph{and} a smaller $\log\N$ to divide by---the two effects
compound, which is why the frontier steepens so sharply below $130$
digits.}\label{tab:frontier}
\begin{tabular}{rrrr}
\toprule
Digits of $\N$ & Moduli & $|U|$ & $E$\\
\midrule
$195$ & $89$ & $692$ & $16.8$\\
$170$ & $79$ & $754$ & $20.6$\\
$150$ & $72$ & $802$ & $24.4$\\
$140$ & $68$ & $830$ & $26.8$\\
$130$ & $64$ & $862$ & $29.5$\\
$120$ & $60$ & $899$ & $32.8$\\
$100$ & $52$ & $973$ & $41.3$\\
\bottomrule
\end{tabular}
\end{table}

Table~\ref{tab:ladder} does the complementary experiment: fix the digit budget
at $199$ and raise the target. This is the budget game's own ladder, and it is
the table that governs Chapter~\ref{ch:outlook}'s planning.

\begin{table}[h]
\centering\small
\caption{The rung ladder at $199$ digits: exponent and candidate count as the
target $Q$ rises, using the best cover we can build at each target ($88$ moduli,
$\mathrm{boost}\approx11$). Fleet-days assume the measured
$7.9\times10^{5}$ candidates/s of a two-session GPU fleet.}\label{tab:ladder}
\begin{tabular}{rrrrr}
\toprule
$Q$ & $|U|$ & $E$ & $e^{E}$ candidates & fleet time\\
\midrule
$122\,000$ & $867$ & $20.83$ & $1.1\cdot10^{9}$ & $23$ minutes\\
$135\,000$ & $1\,004$ & $24.12$ & $3.0\cdot10^{10}$ & $11$ hours\\
$150\,000$ & $1\,138$ & $27.34$ & $7.5\cdot10^{11}$ & $11$ days\\
$175\,000$ & $1\,335$ & $32.07$ & $8.5\cdot10^{13}$ & $3.4$ years\\
$200\,000$ & $1\,569$ & $37.69$ & $2.4\cdot10^{16}$ & $950$ years\\
\bottomrule
\end{tabular}
\end{table}

The spacing is geometric---each rung costs $25$ to $280$ times the last---which
has a pleasant strategic consequence exploited in Chapter~\ref{ch:outlook}:
climbing \emph{every} rung costs only about $4\%$ more than jumping straight to
the top, and banks a certified record at each step.

\section{Cover size is an equilibrium, not a budget}\label{sec:equilibrium}

A natural mental model is that the cover is limited by the digit budget: use as
many moduli as $\log M\le\log\N$ allows. That model is wrong at record scale,
and correcting it explains why the frontier is so flat.

Consider adding one more modulus $r=461$ to the $88$-modulus cover behind our
budget-game work. It kills about $3.5$ residual offsets, worth
$-3.5\times\mathrm{boost}/\log\N=-0.084$ nats. But it also multiplies
$\mathrm{boost}$ by $461/460$, which raises the primality odds of every
\emph{surviving} residual complement, worth $+0.082$ nats. The net effect is
$+0.00$ nats: the cover is at a stationary point.

\begin{observation}[Boost--kills equilibrium]\label{obs:equilibrium}
At $B=10^{200}$ the cover behind our budget-game work spends $\log M\approx437$
of the $\approx461$ available nats, so the digit budget is not tight; the cover
stops growing because $dE=0$, not because it has run out of room. Adding moduli
past the equilibrium is neutral to slightly harmful.
\end{observation}

This is why leftover digit budget buys nothing, why the frontier in
Table~\ref{tab:frontier} is smooth rather than kinked, and why---as we will
see---only two things can actually move $E$: better rounding mathematics, or
more compute.

\section{The linear-programming relaxation is vacuous}\label{sec:lp}

The covering problem is a set-cover: choose one residue class per modulus so as
to cover as many prime offsets as possible. Set-cover problems come with a
standard relaxation---allow fractional choices---whose optimum lower-bounds the
integral one and whose dual prices certify optimality. A natural research
programme, proposed in our own earlier planning documents, was to solve that
relaxation and publish the duals as a certificate of how close greedy is to
optimal.

We solved it. The result kills the programme, and in doing so reframes the whole
subject.

\begin{proposition}[LP vacuousness]\label{prop:lp}
Let $\mathcal R$ be the $88$ moduli $3,5,\dots,457$ used at record scale and let
the universe be the odd primes below $Q$. In the relaxation where each modulus
$r$ may distribute a unit of weight over its residue classes and a prime is
covered to the extent that weight sits on its class, the optimum covers
$100.00\%$ of the primes---at both $Q=122\,000$ and $Q=200\,000$.
\end{proposition}
\begin{proof}[Verification]
Solved with HiGHS in about two seconds at each target. The reason is a mass
count: modulus $r$ covers a $1/(r-1)$ fraction of the odd primes in expectation,
and $\sum_{r\in\mathcal R}1/(r-1)\approx1.86$. With $186\%$ of the required
covering mass available, a fractional assignment can spread it to leave nothing
uncovered.
\end{proof}

So the LP dual certifies nothing whatsoever about integral covers, and the
$\S7.4$-style ``optimality certificate'' programme is dead as posed. What
replaces it is a much better statement:

\begin{corollary}[The exponent is a rounding gap]\label{cor:rounding}
Every nat of the failure exponent is the cost of integrality. Concretely, at
$Q=122\,000$:
\[
\underbrace{15.6\%}_{\text{naive rounding}}\;\longrightarrow\;
\underbrace{7.56\%}_{\text{adaptive greedy}}\;\longrightarrow\;
\underbrace{0\%}_{\text{fractional optimum}}
\]
uncovered primes, where the first figure is $e^{-1.86}$, the density left by
independent rounding at the available mass. Each absolute percentage point of
residual removed is worth about $2.7$ nats---a factor of $15$ in compute.
\end{corollary}

Greedy therefore sits almost exactly halfway, in halvings, between doing nothing
clever and perfection. The rest of this chapter asks whether the remaining half
can be taken.

\section{Attacking the rounding gap}\label{sec:door1}

We fix a concrete benchmark: the $88$ moduli above, the $11\,474$ odd primes
below $Q=122\,000$, and the greedy/coordinate-descent incumbent with
$|U|=867$. One prime saved is worth $0.0240$ nats. The question is how much of
the gap to $0$ is reachable.

\subsection{How much is there to win?}

Before optimising, it is worth knowing the size of the prize. Sampling the space
of assignments uniformly (choosing each $a_r$ at random) gives
$|U|$ with mean $\mu=1377$ and standard deviation $\sigma=21.0$. Two things are
notable. First, $\sigma$ is \emph{below} the independence prediction of $34.8$:
the choices are negatively correlated, so the deep tail is thinner than a naive
model suggests. Second, the design space has entropy
$\ln\prod_r(r-1)=435.7$ nats, so a first-moment (annealed) calculation puts the
best assignment in the ensemble at roughly
\[
u^{*}\;\approx\;\mu-\sqrt{2\ln\Omega}\,\sigma\;\approx\;758 ,
\]
against greedy's $867$. Greedy already sits at $z=24.3$ of the $29.5$ standard
deviations theoretically available. The honest prize is therefore about
$2.6$ nats---a factor of $14$---and only if the annealed threshold is
achievable, which for a clustered landscape it typically is not.

\subsection{Exact local optimisation}

We attacked the benchmark with large-neighbourhood search: freeze all but a
window of $7$--$15$ moduli and re-solve that window \emph{exactly} with
CP-SAT~\cite{CPSAT} over the primes left uncovered by the frozen part, then
rotate windows. Windows were chosen three ways---by tightness (the moduli whose
current class beats its best alternative by the smallest margin), randomly among
tight moduli, and uniformly---and the incumbent was enforced as a hard floor so
the solver could only match or improve.

Across $41$ windows, with time caps up to $300$ seconds, the incumbent was
matched every time and never beaten. To turn ``not beaten'' into ``provably
optimal at this scale'' we added an LP certificate per window: solve the
window's own relaxation, and if $\lfloor\mathrm{LP}\rfloor$ equals the
incumbent's coverage then no integral assignment of that window can do better.
This certified $21$ of $22$ windows where it was computed. The single exception
was a window of the smallest moduli---where, exactly as
Proposition~\ref{prop:lp} predicts, the LP is loose.

We then went further and solved the \emph{entire large-modulus endgame} at once:
first the $30$ largest moduli simultaneously, then the $44$ largest---half the
cover in a single exact subproblem. Both leave $|U|=867$.

\subsection{Global methods}

Local exactness having failed, we tried methods that can cross basins.

\emph{Belief propagation with decimation.} We wrote the assignment problem as a
factor graph---one variable per modulus with domain $\mathbb{Z}/r$, one factor
per prime contributing weight $e^{-\beta}$ when uncovered---and ran sum-product
message passing, then repeatedly froze the most polarised modulus and re-ran
(the standard decimation scheme~\cite{MezardMontanari}). Validation first: on a
small instance where exhaustive enumeration is possible, BP tracks the exact
$E[|U|]$ to $0.7\%$ in the soft regime, and BP-guided decimation finds the exact
optimum. On the real instance it returns $874$ at $\beta=1.5$ (and worse at
higher $\beta$, where early freezes over-commit).

That is seven primes short of greedy, and it is the second-best construction we
know for this instance---but it is not an improvement.

\emph{Basin structure.} The most informative experiment was the cheapest. Take a
solution produced by random-start coordinate ascent, which lands at
$|U|\approx968$, and apply the same exact-window machinery to it: twenty
windows recover \emph{none} of its $101$-prime gap to greedy. Independent ascent
solutions agree with each other on only $2.4\%$ of their class choices---chance
level---so there is no backbone to exploit, and crossover-style population
methods have nothing to graft.

\begin{observation}[The landscape]\label{obs:landscape}
The assignment space is a field of mutually distant, locally-exact-stable
basins spanning at least $|U|\in[867,980]$. Exact local search cannot move
between them at all; what distinguishes greedy is not the neighbourhood it
explores but the \emph{sequentially adaptive} way it constructs a solution,
which lands it about $100$ primes deeper than a typical basin.
\end{observation}

\subsection{Verdict, and one methodological caution}

The greedy plateau survived every method we brought: $41$ exact windows (mostly
LP-certified), a half-cover exact block, BP decimation, and basin probing. We
state the conclusion carefully, because the temptation to over-claim here is
strong.

\begin{observation}[Door-1 verdict]\label{obs:door1}
At $Q=122\,000$ with these moduli, $|U|=867$ is at or very near the quenched
optimum. The annealed estimate of $758$ is an \emph{ensemble-existence}
heuristic, not a constructive promise, and the clustering visible in
Observation~\ref{obs:landscape} is exactly the situation in which annealed and
quenched optima diverge. The realistic prize is therefore $0$ to $2.6$ nats,
probably nearer the bottom of that range.
\end{observation}

The caution concerns mean-field methods. Naive belief propagation, run to low
temperature, predicts $E[|U|]\to35$ at $\beta=4$---heading for the fractional
fantasy of Proposition~\ref{prop:lp} rather than the achievable $867$. The
replica-symmetric approximation inherits precisely the LP's blindness to
integrality. Anyone using message passing to estimate a threshold here needs the
one-step replica-symmetry-breaking treatment with complexity bookkeeping; we
flag this as the honest remaining theory item rather than quoting a number we do
not trust.

\section{What the frontier costs in money}

Putting Table~\ref{tab:frontier} together with measured throughput gives the
price list in Table~\ref{tab:menu}. It is worth stating in these terms because
the strategic question in Chapter~\ref{ch:outlook} is precisely which of these
rungs is worth buying.

\begin{table}[h]
\centering\small
\caption{Cost of the threshold game at $Q=10^{5}$ by digit target, at
$8.5\times10^{6}$ tests/s per A100 and cloud pricing current at the time of
writing.}\label{tab:menu}
\begin{tabular}{rrrr}
\toprule
Digits & $E$ & A100-years & approx.\ cost\\
\midrule
$140$ & $28.1$ & $0.10$ & \$1\,000\\
$135$ & $29.7$ & $0.49$ & \$5\,000\\
$130$ & $31.2$ & $2.3$ & \$24\,000\\
$125$ & $32.9$ & $12$ & \$125\,000\\
$120$ & $34.7$ & $73$ & \$770\,000\\
$100$ & $44.1$ & $8.9\cdot10^{5}$ & \$9.4\,billion\\
\bottomrule
\end{tabular}
\end{table}

The table makes the shape of the problem vivid. The frontier is fundable down to
about $130$ digits and institutional down to $120$; sub-$100$ is seven orders of
magnitude beyond plausible compute. Whether that last wall can be attacked by
mathematics rather than money is the subject of the next chapter.

\chapter{Sub-100 digits: a reduction to two hard problems}\label{ch:barriers}

Chapter~\ref{ch:frontier} priced the threshold game at $Q=10^{5}$ and found the
$100$-digit rung seven orders of magnitude beyond plausible compute. This
chapter asks whether that wall is real. The answer is that it is real
\emph{within the paradigm}, and that stepping outside the paradigm leads
directly to two well-known hard problems---which is a more useful outcome than
either a proof or a construction, because it says exactly what would have to be
broken.

It is worth being clear at the outset that sub-$100$-digit deserts \emph{exist}
in abundance. With no cover at all, the failure exponent at $B=10^{100}$ is
$E_{\text{random}}\approx110$, so the density of qualifying integers is
$e^{-110}$ and the expected count below $10^{100}$ is about
$10^{100}\cdot e^{-110}\approx10^{52}$. Nothing here is an existence
obstruction. The entire difficulty is constructive.

\section{Two regimes}

Write $A=\log M$ for the cover mass in nats and $\log B$ for the digit budget
($\log B\approx230$ at $100$ digits). Covers split into two regimes according to
whether their modulus fits under the budget.

\paragraph{Regime I: $A\le\log B$.} The CRT progression itself supplies
$e^{\log B-A}$ candidates, and the search costs $e^{E}$. Optimising at
$Q=10^{5}$, $B=10^{100}$ gives a best cover of $44$ moduli with
$M\approx10^{79}$, $|U|=1060$, $\mathrm{boost}=9.6$, hence
\[
E\;=\;44.1,\qquad e^{E}\approx1.4\times10^{19}\ \text{candidates}.
\]
At the record engine's measured throughput of about $5\times10^{5}$ candidates
per second this is roughly $9\times10^{5}$ GPU-years.

\begin{remark}[A units error, corrected]\label{rem:units}
An earlier version of this calculation reported $5\times10^{4}$ GPU-years, an
error of a factor of $18$: it divided the candidate count by the \emph{test}
rate ($8.5\times10^{6}$ Fermat tests/s) rather than the \emph{candidate} rate,
forgetting that a candidate costs about $17$ tests. The error was caught by
cross-checking against an independent study of the same frontier, which measured
a realisable $E=43.66$ at $100$ digits---within $0.4$ nats of the figure above,
and irreconcilable with the faster cost. We record the correction here because
the discipline that caught it (always cross-check a headline cost against an
independently derived one) is worth more than the number.
\end{remark}

\paragraph{Regime II: $A>\log B$ (oversized covers).} If the modulus is allowed
to exceed the budget, coverage becomes excellent. All $9\,592$ primes below
$10^{5}$ can be covered by $355$ moduli of total mass $1\,019$ digits; allowing
each modulus to use $L=16$ residue classes, $420$ moduli leave only $|U|=165$,
giving
\[
E\;=\;10.2,\qquad e^{E}\approx27\,000\ \text{candidates}.
\]
Twenty-seven thousand candidates is nothing. If we could produce such a design
\emph{together with} a representative below $10^{100}$, sub-$100$ would fall
this afternoon.

Solutions exist: the design space (which moduli, which classes) carries about
$2\,698$ nats of entropy against $\log(M/B)=2\,630$ nats of ``distance to
squeeze'', leaving a surplus of about $+57$ nats and hence roughly $e^{57}$
valid (design, $\N$) pairs below $10^{100}$. But the \emph{density} of designs
whose representative lands below $B$ is $e^{-2630}$. Enumeration is out by some
$1\,100$ orders of magnitude.

\begin{observation}
Regime II is where sub-$100$ becomes easy \emph{if and only if} small
representatives can be found directly. That is the whole problem.
\end{observation}

\section{The small-representative problem is modular subset-sum}

CRT reconstruction is additive. Writing the standard idempotent basis,
\[
\N\;\equiv\;\sum_{r\in S}a_r\cdot\frac{M}{r}\cdot\Bigl(\bigl(\tfrac{M}{r}\bigr)^{-1}\bmod r\Bigr)
\pmod M ,
\]
so ``find a design whose representative is below $B$'' is exactly:
\begin{quote}\itshape
minimise $\bigl|\sum_r c_r \bmod M\bigr|$ over choices $c_r\in C_r$ with
$|C_r|=L$, targeting the window $[0,B)$.
\end{quote}
This is an inhomogeneous \emph{modular subset-sum problem with per-position
choices}. Its density---the ratio of available design entropy to the window
squeeze---is
\[
d\;=\;\frac{2\,698}{2\,630}\;\approx\;1.026 .
\]

That is the worst possible value. Low-density instances ($d<0.94$) fall to
lattice reduction~\cite{LagariasOdlyzko,Coster}; very high density instances
fall to birthday and $k$-tree methods~\cite{Wagner}. Density $\approx1$ is
precisely the regime where neither works and where the presumed hardness of
subset-sum is concentrated. Concretely, our solution set has size about
$2^{82}$ inside a design space of about $2^{3893}$; meet-in-the-middle gives
$2^{1897}$ and Wagner's $k$-tree with our list sizes still needs $2^{938}$.
Both are astronomically short.

\section{The decoding shortcut, and why it fails}

There is one route that would bypass the subset-sum problem entirely. Finding a
small integer consistent with many congruences is a decoding problem for the
Chinese Remainder code, and CRT codes have list-decoding
algorithms~\cite{GSS,Goldreich}. If the decoder's radius were large enough, we
could simply ask it for a small $\N$ agreeing with a chosen cover.

The relevant statement is Guruswami--Sahai--Sudan~\cite{GSS}, Theorem 3: for a
CRT code with moduli $p_i$, message bound $K$, and any non-negative integers
$\ell$ and $z_i$, one finds in time $\mathrm{poly}(n,\log N,\ell,\sum z_i)$ all
codewords $m$ with
\[
\sum_i \alpha_i z_i\log p_i\;>\;\log(\ell+1)+\frac{\ell}{2}\log\frac{K}{2}
+\frac{1}{\ell+1}\sum_i\binom{z_i+1}{2}\log p_i ,
\]
where $\alpha_i=1$ exactly when $m\equiv r_i\pmod{p_i}$. Optimising $z$ and
$\ell$ under uniform weights, with $L$ allowed classes per position, gives an
agreement-mass radius
\begin{equation}\label{eq:gss}
A\;>\;\sqrt{L\cdot\log B\cdot P},\qquad P=\sum_{\text{pool}}\log p_i .
\end{equation}
We verified~\eqref{eq:gss} numerically against the exact finite-$\ell$ optimum:
the ratio is $1.001$--$1.004$ across pools from $5\times10^{3}$ to $10^{5}$ and
$L\in\{1,2\}$.

Now the trap closes. The radius scales with the mass $P$ of the \emph{whole
pool}, which must be far larger than the cover itself in order to supply the
subset entropy that makes small representatives exist at all. Numerically, at
pool bound $60\,000$ with $L=16$: the pool mass is $59\,816$ nats, the decoding
radius is $14\,845$ nats, and our cover mass is $2\,860$ nats. The decoder needs
agreement $5.2\times$ larger than anything we can supply. Shrinking the pool
lowers the radius but destroys the entropy that makes solutions exist.

\begin{theorem}[Barrier; heuristic, first-moment]\label{thm:barrier}
For CRT-cover constructions of even $\N<B$ with $\g(\N)>Q$, no parameter choice
simultaneously satisfies
\begin{enumerate}
\item[(a)] $A<\log B+\log\binom{n}{s}+s\log L$ \quad (representatives exist), and
\item[(b)] $A>\sqrt{L\cdot\log B\cdot P}$ \quad (GSS-decodable).
\end{enumerate}
Hence sub-$100$ via oversized covers requires either list-recovery of CRT codes
\emph{beyond} the Johnson-type radius---an open problem in coding theory---or a
subset-sum algorithm at density $\approx1$.
\end{theorem}

\begin{table}[h]
\centering\small
\caption{The barrier, evaluated. ``Exists'' is the largest cover mass for which
representatives below $B$ still exist; ``decodes'' is the GSS radius. The ratio
never approaches $1$.}\label{tab:barrier}
\begin{tabular}{rrrrrr}
\toprule
pool bound & $L$ & pool mass & $A$ exists & $A$ decodes & ratio\\
\midrule
$1\,000$ & $1$ & $956$ & $339$ & $469$ & $0.72$\\
$60\,000$ & $1$ & $59\,816$ & $570$ & $3\,711$ & $0.15$\\
$60\,000$ & $16$ & $59\,816$ & $1\,305$ & $14\,845$ & $0.09$\\
$10^{6}$ & $16$ & $998\,483$ & $1\,697$ & $60\,651$ & $0.03$\\
\bottomrule
\end{tabular}
\end{table}

\subsection{A retracted claim, and what it teaches}\label{sec:retraction}

An earlier version of this analysis reported an \emph{open window}: a range of
pool sizes where, apparently, representatives existed and the decoder still
worked. That claim was wrong and is withdrawn. The error is instructive enough
to record.

The two conditions of Theorem~\ref{thm:barrier} both depend on an agreement
size $s$---how many positions the sought codeword matches. The retracted
analysis computed the decoding radius at one value of $s$ and the existence
bound at a \emph{different} value (the one maximising existence, near $2n/3$),
and compared the two numbers. Requiring both conditions at the \emph{same} $s$
closes the window at every pool size. The lesson generalises well beyond this
problem: \emph{a feasibility window assembled from two separately optimised
bounds is worthless; couple the parameters first, then optimise once.}

We also had the barrier re-derived independently, by a route that shares no
code with ours---a Legendre-transform treatment of the same optimisation---which
reproduces the $14\,845$-nat radius and finds the gap closed at every pool size,
with the best case still $28\%$ short. Independent confirmation matters
especially when the first version of a result was wrong.

\section{Does lattice reduction beat the radius in practice?}

Decoding radii are worst-case statements. Our instances are highly structured
and, as computed above, admit roughly $e^{57}$ solutions---so it is entirely
reasonable to hope that lattice reduction succeeds far beyond the proved radius,
as it often does elsewhere. We tested this directly rather than assuming either
way.

We built a small CRT-decoding testbed: plant a known small solution in a
random instance at a controlled agreement mass, then run
LLL~\cite{LLL} on the standard Howgrave-Graham-style lattice and record whether
the planted solution (or any other) is recovered. Sweeping the agreement mass
across the Johnson-type radius $A_J$ gives a sharp threshold:

\begin{itemize}
\item at $\ell=30$, recovery succeeds above $1.028\,A_J$ and fails below;
\item at $\ell=45$ the threshold tightens to about $1.02\,A_J$, converging to
the theoretical radius from above;
\item below the threshold the failure is \emph{total}: across $390$ decodes we
recovered zero non-planted roots, and no partial or near-solutions.
\end{itemize}

The last point is the one that matters. Solution abundance buys nothing: even
with $10^{7}$ alternative solutions available in the instance, the lattice
returns nothing at all below the radius. The Johnson-type radius behaves here
like a genuine algorithmic boundary rather than a proof artifact, which is what
makes Theorem~\ref{thm:barrier} a barrier and not merely a gap in our knowledge.

\section{Where this leaves sub-100}

\begin{enumerate}
\item \textbf{Regime I is the only currently viable route}, and it is out of
reach: $E=44.1$, about $9\times10^{5}$ GPU-years, seven orders of magnitude
beyond plausible compute. The fundable frontier is $130$--$140$ digits
(Table~\ref{tab:menu}).

\item \textbf{Two well-posed open problems now sit under the record}, either of
which would collapse it: (i) list-recovery of CRT codes beyond the Johnson-type
radius for structured, solution-abundant instances; (ii) modular subset-sum with
per-position choices at density $\approx1$ where $2^{82}$ solutions exist. Both
are of independent interest, which is the right way for a chapter like this to
end.

\item \textbf{The one place the barrier does not apply} is a construction that
certifies compositeness by something other than a congruence divisor. Every
argument in this chapter accounts entropy in units of $\log r$ per excluded
offset; a mechanism that escapes that accounting escapes the barrier. Whether
any such mechanism exists is the subject of Chapter~\ref{ch:closure}, and the
answer, within a precisely defined language, is no.
\end{enumerate}

\chapter{A closure theorem for the covering paradigm}\label{ch:closure}

Every record in this thesis, and every record we know of in this problem, is
built the same way: choose $\N$ so that each small prime offset $q$ has
$\N-q$ divisible by some small prime fixed in advance. Chapter~\ref{ch:barriers}
showed that pushing this paradigm to sub-$100$ digits runs into two hard
problems. The obvious next question is whether the paradigm itself can be
escaped---whether some other mechanism could certify $\N-q$ composite more
cheaply than a congruence.

This chapter argues that it cannot, within a language precise enough to make the
claim checkable. The result is not a theorem in the sense of a completed proof
in a journal; it is a theorem in the sense of a stated claim, with a proof
skeleton whose cases are individually established, two named conditionals, and
an adversarial falsification campaign that failed to break it. We are explicit
throughout about which is which.

\section{The empirical precursor}\label{sec:wp7}

Before formalising anything, we ran five bounded probes into candidate
mechanisms, each with a predeclared quantitative gate. The metric is
\emph{compression}: offsets certified per nat of design entropy, benchmarked
against covering (which manages roughly $4\,400$ offsets per nat at $r=3$,
falling to about $0.4$ per nat at $r\approx1\,200$).

\begin{description}
\item[P1. Prime-power moduli.] Upgrading a modulus $p$ to $p^{2}$ costs an extra
$\log p$ nats and yields \emph{zero} new kills: the class $q\equiv a\pmod{p^{2}}$
is a subset of the class mod $p$ that is already killed (at $p=7$ we count
$227\subset1\,613$). Strictly dominated.

\item[P2. Polynomial value-set identities.] Take $\N=m^{k}-c$, so that
$\N-q=m^{k}-d^{k}$ is algebraically composite whenever $q=d^{k}-c$. The best
case is $k=2$: $c=398$ kills $113$ offsets below $10^{5}$. But confining $\N$ to
a $k$-th-power family surrenders $(1-1/k)\log B$ nats---$115$ nats at
$B=10^{100}$---and the value set below $Q$ has at most $O(Q^{1/2})$ members. The
mechanism cannot scale past a few hundred offsets \emph{at any price}.

\item[P3. Sierpi\'nski--Riesel exponent families.] For $\N_n=k\cdot2^{n}+c$ the
residue mod $p$ cycles with period $\mathrm{ord}_p(2)$. For any fixed candidate
the induced constraint is one congruence class: ordinary covering, no entropy
gain. (The exponent structure does correlate residues \emph{across} the family,
which is a legitimate search-throughput trick---noted for the engine, not for
the entropy budget.)

\item[P4. Norm forms and splitting conditions.] ``$\N-q$ is a norm in $K$'' does
not certify compositeness---rational primes with the right splitting behaviour
are norms---and any usable splitting condition on $q$ is a congruence condition
modulo discriminant-related moduli. Covering in Galois-theoretic clothing.

\item[P5. Residual bias.] Could $\N$ be chosen so that residual complements are
composite \emph{more often} than the boost predicts? The banked campaigns answer
directly: across three covers and $1.2\times10^{9}$ candidates,
$|\Ehat-E|\le0.1$. Any unexploited compositeness structure is worth at most
$0.1$ nats.
\end{description}

All five terminated at their gates. That is suggestive but not conclusive: five
probes is five probes. The rest of this chapter turns the suggestion into a
claim with a proof structure, a formal language, and an adversary.

\section{The formal setting}

\begin{definition}[Construction scheme]\label{def:scheme}
Fix a bound $B$ and target $Q$. A \emph{construction scheme} is a pair
$(F,C)$:
\begin{itemize}
\item a \emph{candidate family} $F$: a parameter set $\Theta\subseteq\mathbb Z^{m}$
with a poly-size description and an evaluation map
$\N:\Theta\to2\mathbb Z\cap[1,B)$ computable in $\mathrm{poly}(\log B)$ time. Its
\emph{design cost} is $\mathrm{cost}(F)=\log B-\log|\N(\Theta)|$, the nats of
search freedom surrendered;
\item a \emph{certificate map} $C$: on input $(\theta,q)$ with $q<Q$ prime, it
outputs $\bot$ or a divisor witness $D(\theta,q)$ with
$1<D<\N(\theta)-q$ and $D\mid\N(\theta)-q$, verifiable in $\mathrm{poly}(\log B)$
time.
\end{itemize}
The certificate is \emph{planted} if the divisibility holds as an algebraic
identity over the family (checkable symbolically), and \emph{extracted} if $C$
discovers the divisor per instance. Write $K(\theta)$ for the certified offsets
and $U(\theta)$ for the residual ones.
\end{definition}

The \emph{effective exponent} $E_{\text{eff}}(F,C)$ is the logarithm of the
expected work, in candidates, to find a $\theta$ all of whose residual
complements are composite. For the classical cover, $E_{\text{eff}}=E$ of
\eqref{eq:E}, calibrated to $\pm0.1$ nats in \S\ref{sec:calibration}.

\begin{theorem}[Divisor-Paradigm Closure; target statement]\label{thm:closure}
Let $(F,C)$ be a construction scheme whose certificate map is planted and
expressible in the schema language $L$ of \S\ref{sec:grammar}. Assume:
\begin{itemize}
\item[(H1)] Hardy--Littlewood-type heuristics for shifted primes: conditional on
all divisibility data, the primality indicators of distinct complements are
independent with the Mertens-adjusted rates;
\item[(H2)] no compositeness test with constant soundness on random rough inputs
runs in $o(\text{one modular exponentiation})$.
\end{itemize}
Then $E_{\text{eff}}(F,C)\ge E^{*}(Q,B)-o(1)$, where $E^{*}$ is the covering
frontier. Moreover, extracted certificates on any window-dense family of targets
are factoring-equivalent, and certificate-free detection is already performed at
its known-optimal cost by the standard engine.
\end{theorem}

Informally: \emph{within $L$, nothing beats covering; outside planted-$L$, a
scheme must either factor or test, and both are at known cost floors.}

\section{The three lemmas}

\subsection{Lemma 2 first: correlations are covering in disguise}

We take the lemmas out of order because the second is the one that is actually
proved, and it does a surprising amount of work.

\begin{proposition}\label{prop:shared}
Let $\N$ be even and $q_1\ne q_2$ odd primes below $Q$. If $d>1$ divides both
$\N-q_1$ and $\N-q_2$, then $d\mid q_1-q_2$, and hence every prime factor of $d$
is smaller than $Q$.
\end{proposition}
\begin{proof}
Subtract: $d$ divides $(\N-q_1)-(\N-q_2)=q_2-q_1$, and $0<|q_1-q_2|<Q$.
\end{proof}

Trivial to prove, and yet it settles an entire class of speculative mechanisms.
Any scheme that hopes to make the residual complements \emph{jointly} luckier---%
choosing $\N$ so that its leftovers conspire to be composite together---must act
through a divisor shared by two complements, and Proposition~\ref{prop:shared}
says every such divisor is supported on primes below $Q$. But divisibility by
small primes is exactly what the cover fixes and the sieve harvests. There is
nothing left to exploit.

\begin{lemma}[Correlation exhaustion]\label{lem:correlation}
Let $\mathcal F_{<Q}$ be the information generated by the events
$\{s\mid\N-q\}$ for primes $s,q<Q$. Then every two-complement divisor event is
$\mathcal F_{<Q}$-measurable; conditional on $\mathcal F_{<Q}$ (and on the
sieve's data up to its bound), (H1) makes the residual primality indicators
independent, so the success probability factorises as $\prod_q(1-p_q)$ and
$E_{\text{eff}}$ decomposes as design cost plus $\sum_q p_q$---which is the
covering objective. Empirically, $|\Ehat-E|\le0.1$ over $1.2\times10^{9}$
candidates bounds any violation.
\end{lemma}

\subsection{Lemma 1: entropy accounting for planted certificates}

\begin{lemma}[Entropy accounting]\label{lem:entropy}
Let $(F,C)$ be planted. For each mechanism define its \emph{certified sub-image
density} $\delta_{\text{eff}}$: the density, within a generic window of width
$Q$ at height $B$, of integers for which the scheme can exhibit its planted
witness using poly-range parameters. Then the amortised design cost per
certified offset is at least $\log(1/\delta_{\text{eff}})-O(\log\log B)$, and
for every family in $L$ this is at least the covering cost.
\end{lemma}

The proof is by cases, and the cases are where the content is. We treat them in
\S\ref{sec:cases} after fixing the language, because the arguments are
family-specific and quantitative rather than uniform.

\subsection{Lemma 3: the detector floor, as an open problem}

Hypothesis (H2) deserves to be stated on its own, because it is the one
genuinely open ingredient and it is interesting independently of Goldbach
deserts.

\begin{conjecture}[H2, sub-modexp compositeness detection]\label{conj:h2}
Let $n$ be uniformly random among integers in $[X,2X]$ with no prime factor
below $B_0=(\log X)^{O(1)}$. Does there exist a randomised algorithm deciding
compositeness of $n$ with soundness $\ge2/3$ in time
$o(M(\log X)\cdot\log X)$---asymptotically cheaper than one modular
exponentiation?
\end{conjecture}

Known art: trial division and batch-gcd methods~\cite{Bernstein} certify exactly
the small-divisor mass and have soundness $o(1)$ on rough inputs; Jacobi symbols
are equidistributed on rough composites; Fermat, Euler, Miller--Rabin and
Frobenius tests all cost $\Theta(M\log X)$. We are aware of no sub-modexp result
in either direction, and no lower bound in any realistic model.

The downstream consumer is quantified: a detector at cost $c$ and recall $\rho$
improves throughput by $1/(c+(1-\rho))$, so $c=0.1$, $\rho=0.9$ is a $5\times$
speedup on every search in this thesis---and, since record difficulty is
logarithmic in compute, converts directly into records. The same primitive
dominates GIMPS and PrimeGrid~\cite{PrimeGrid}, so a positive answer would be
felt well beyond this problem.

\section{The schema language \texorpdfstring{$L$}{L}}\label{sec:grammar}

To make Theorem~\ref{thm:closure} checkable we fix a concrete grammar of
certificate mechanisms. Each schema compiles to three things: an enumerator of
the offsets it certifies, a symbolic divisor witness, and a design-cost term.
\[
\begin{array}{ll}
L::= & \mathrm{Congruence}(d,a)\quad \N\equiv a\ (\mathrm{mod}\ d)\\
   |\ & \mathrm{PolyImage}(f\in\mathbb Z[m],\deg\le4)\quad \N-q=A(m)B(m)\\
   |\ & \mathrm{ExpFamily}(\N=j\cdot a^{n}+c)\quad\text{per-offset order conditions}\\
   |\ & \mathrm{MultiPoly}(A,B\in\mathbb Z[u_1..u_v],\deg\le2,\ v\le3)\\
   |\ & \mathrm{NormForm}(K,\ \mathrm{disc}\le10^{4})\quad \N-q=\mathrm{Norm}_K(\alpha)\\
   |\ & \mathrm{Cyclotomic}(\Phi_n,\ n\le200)\\
   |\ & \mathrm{Compose}(L,L)\quad\text{depth}\le2
\end{array}
\]
Two exclusions are deliberate. Mechanisms whose certificate requires a
primality decision to define membership are circular and excluded; and
per-instance divisor discovery is \emph{extraction}, handled by the second horn
of the theorem rather than by $L$.

Two thresholds make ``beats covering'' precise. A schema must certify more than
$\mathrm{CEIL}=3\sqrt{Q}\approx1\,048$ distinct prime offsets---placing it above
the $\sqrt{Q}$ ceiling that limits identity families---and it must do so at a
marginal cost below covering's. The right benchmark for the second condition is
subtle and we got it wrong at first (\S\ref{sec:referee}); the correct one is
covering's \emph{average}:
\[
\frac{437\ \text{nats}}{10\,607\ \text{offsets}}\;=\;0.041\ \text{nats per offset},
\]
not its marginal endgame cost of about $1.01$. Covering wins on average, not at
the margin, because one modulus spends $\log r$ nats and certifies
$\pi(Q)/(r-1)$ offsets \emph{at once}, whereas every known algebraic mechanism
buys $O(1)$ offsets per unit of entropy. Together with the design-entropy budget
$\log B\approx461$, the two conditions imply that a bulk mechanism needs
$\mathrm{cost}<\log B/\mathrm{CEIL}=0.44$ nats per offset.

\section{Closing the families}\label{sec:cases}

\subsection{Univariate identity families: a ceiling, not a price}

A degree-$d$ form takes at most $Q^{1/d}$ values below $Q$, so it can certify at
most that many offsets \emph{regardless of budget}: $349$ at $d=2$, $50$ at
$d=3$, $19$ at $d=4$, $7$ at $d=6$. Since $\mathrm{CEIL}=1\,048$, every
$\mathrm{PolyImage}$ with $d\ge2$ is closed at any price, and $d=1$ is ordinary
covering. Cyclotomic schemas fall inside the same bound because $\varphi(n)\ge2$
for every $n\ge3$. This is the P2 probe generalised, and it retires two families
outright.

\subsection{ExpFamily: supply, not price}

$\mathrm{ExpFamily}$ is the strongest non-covering family in the grammar and
needs a quantitative argument. For $\N=j a^{n}+c$, a prime $s$ certifies the
offset $q$ when $(q-c)j^{-1}$ lies in the subgroup $\langle a\rangle$ mod $s$,
and the resulting condition confines the exponent $n$ to one class modulo
$\mathrm{ord}_s(a)$: a cost of $\log\mathrm{ord}_s(a)$ nats \emph{per offset}.

Two facts then close it. First, a structural floor: $s\mid a^{d}-1$ forces
$a^{d}>s$, so $\mathrm{cost}=\log d>\log\log_a s$---cheap certificates require
small moduli, which is exactly where covering is cheapest. Exhaustively over all
$s<10^{5}$ the cheapest is $\log 5=1.609$ (at $a=7$, $s=2801$).

Second---and this is the binding constraint---cheap certificates are
\emph{scarce}. A modulus $s$ serves only the offsets landing in its order-$d$
subgroup, a fraction $d/s$ of them, so it supplies about $\pi(Q)d/s$ offsets:
$s=2801$ with $d=5$ supplies about twenty. Solving the resulting purchase
problem exactly---sort every $s<10^{5}$ by $\log\mathrm{ord}$, buy offsets
cheapest-first---gives, for every base $a\in\{2,3,5,7\}$:

\begin{observation}[ExpFamily is over budget]\label{obs:expfamily}
Reaching $1\,048$ certified offsets costs $1\,151$ nats, against a total design
budget of $\log B\approx461$ nats: $2.5\times$ over. The most that any
budget-feasible $\mathrm{ExpFamily}$ design can certify is $419$ offsets.
\end{observation}

\subsection{MultiPoly: clusters, isolation, and orbits}

Multivariate identities are the interesting case, and the one where we initially
expected the theorem to break. Products of two binary quadratic forms have
\emph{window-dense} value sets---about $9\,300$ of the $17\,983$ prime offsets
below $Q=200\,000$ are representable---which blows straight past the univariate
$\sqrt Q$ ceiling. If those representations could be exhibited, the family would
be a genuine threat.

\paragraph{The cluster bound.} Let $G=A\cdot B$ be homogeneous of total degree
$d$ in $v$ variables, with value-set density exponent $\alpha=\min(1,v/d)$
(so the number of distinct values below $X$ is about $X^{\alpha}$). A
first-moment count of $k$-element clusters of $G$-values inside a width-$Q$
window below $X$ gives $X^{\alpha}(X^{\alpha-1}Q)^{k-1}$, hence clusters exist
only up to
\[
k_{\max}(\alpha)\;=\;1+\Bigl\lfloor
\frac{\alpha\log X}{(1-\alpha)\log X-\log Q}\Bigr\rfloor ,
\]
finite whenever $\alpha<1-\log Q/\log X$. At $X=10^{200}$, $Q=2\times10^{5}$
that threshold is $\alpha^{*}=0.9735$, and every combination with $v<d$ has
$\alpha\le(d-1)/d\le0.9$:

\begin{table}[h]
\centering\small
\caption{Cluster ceilings. A scheme needs about $18\,000$ certified offsets per
candidate; every shared-variable identity family with $v<d$ delivers $O(1)$.}
\label{tab:kmax}
\begin{tabular}{lrr@{\qquad}lrr}
\toprule
$(d_1,d_2,v)$ & $\alpha$ & $k_{\max}$ & $(d_1,d_2,v)$ & $\alpha$ & $k_{\max}$\\
\midrule
$(2,2,2)$ & $0.500$ & $2$ & $(2,2,3)$ & $0.750$ & $4$\\
$(2,3,3)$ & $0.600$ & $2$ & $(2,3,4)$ & $0.800$ & $5$\\
$(2,4,4),(3,3,4)$ & $0.667$ & $3$ & $(2,4,5),(3,3,5)$ & $0.833$ & $6$\\
$(3,4,5)$ & $0.714$ & $3$ & any $v<d$ & $\le0.9$ & $\le13$\\
\bottomrule
\end{tabular}
\end{table}

We checked the density exponent empirically for the one in-grammar case with
$v\ge3$: for $(x^2+y^2+z^2)(x^2+2y^2+3z^2)$, counting distinct values in
width-$10^{4}$ windows at heights $10^{6},10^{8},10^{10}$ gives
$\hat\alpha=0.761$ against the predicted $0.750$.

\paragraph{Above the threshold: isolation.} For $v\ge d$ the value set is
window-dense and existence no longer obstructs, so the closure must come from
\emph{access}. Around any known solution $u_0$, the lattice points $\delta$
keeping $G(u_0+\delta)$ inside a width-$2Q$ window form an ellipsoidal slab with
principal extents $Q X^{-(d-1)/d}$ along the gradient and
$Q^{1/2}X^{-(d-2)/2d}$ along the Hessian directions. At $X=10^{200}$, $Q=2\times
10^{5}$, $d=4$ these are $10^{-145}$ and $10^{-47}$: every extent is far below
$1$, so the only lattice point is $\delta=0$. Window solutions are pairwise
isolated; dense existence is non-constructive for geometric reasons, not
complexity-theoretic ones.

\paragraph{Symmetry cannot help.} One might hope an infinite symmetry group
$\Gamma$ moves between solutions. Three exhaustive cases: if $\Gamma$ preserves
$G$, navigation conserves the value---one offset per orbit. If $\Gamma$ preserves
one factor, it conserves that factor's value, so the certificate always exhibits
the same divisor; if that divisor is small enough to enumerate, the scheme is
sieve-equivalent, and if it is balanced ($\approx X^{1/2}$), we are back to the
isolated sliver. If $\Gamma$ mixes generators from both factors, the reachable
$(\log A,\log B)$ pairs form a coarse lattice with about $(\log X)^{2}\approx
2\times10^{5}$ points against a target window of relative width $10^{-195}$.
Any remaining route must solve $G(u)=m$ for \emph{prescribed} $m$---but a
product-form representation of $m$ \emph{is} a nontrivial factorisation of $m$,
which is the theorem's extraction horn, not a new mechanism.

\subsection{Compose: empty, not expensive}

Finally, depth-2 composition. Our scorer treats a composition's parts as
independent and \emph{sums} their certified offsets---generous---and Compose
still peaked well below threshold. But the generosity conceals a stronger fact:
two nontrivial family constraints on the same $\N$ are contradictory at record
scale. Requiring $\N=m^{2}-c_1$ and $\N=m'^{2}-c_2$ forces
$(m-m')(m+m')=c_1-c_2$, so $m\le|c_1-c_2|$, whereas $\N\sim10^{200}$ needs
$m\sim10^{100}$. For ordinary shifts the solution sets are tiny---$c=398,4686$
admits five solutions, the largest with $m=1073$. The intersection is empty at
scale, so the only viable composition is family-plus-congruence: covering, plus
one family capped at $349$.

\section{The falsification engine}\label{sec:engine-falsify}

A closure claim invites refutation, so we built the refuter ourselves. The
engine compiles a schema to (enumerator, witness, cost), evaluates it against
the live cover context, and applies the filter: a schema \emph{passes} if it
certifies more than $\mathrm{CEIL}$ distinct offsets at below-benchmark cost,
with the design budget enforced.

\paragraph{Tier 0: calibration.} Before trusting the engine we required it to
reproduce the hand-derived numbers of \S\ref{sec:wp7} from primitives. It does:
the best $m^{2}-c$ family below $Q=10^{5}$ is $c=398$ with $113$ killable
offsets, matching P2 exactly, and the apparent discrepancy in the
nats-per-residual-kill figure resolved as a $B=10^{100}$ versus $B=10^{200}$
context difference. Calibration also \emph{corrected} the record: covering's
endgame is about $1.0$ nats per offset at $Q=122\,000$, not the $2.5$ quoted
earlier from a different target.

\paragraph{Tier 1: exhaustive.} We swept $13\,661$ schema instances across all
seven families. \textbf{Zero passes.} The first run reported sixteen, which
turned out to be a filter-semantics leak---congruence-containing schemas being
scored against the escape criterion, when the covering paradigm cannot by
definition escape itself. After hardening, the near-miss board is topped by
$\mathrm{ExpFamily}$ at $3.1$--$4.2$ nats per residual kill, about $3\times$
covering's endgame, exactly as P3 predicted.

\paragraph{Tier 2: generative.} Finally we handed the problem to an adversary:
an LLM-driven evolutionary search~\cite{AlphaEvolve} whose fitness function was
our scorer. The contract was designed to be unfoolable: candidate programs may
\emph{only} emit schema JSON, they run in a locked-down subprocess with no
network, grading happens in a fresh interpreter with the trusted scorer, and a
pass requires hand audit before it counts. A hostile candidate that simply
prints $\texttt{fitness=1.0}$ verifiably scores $0$.

Forty-one evolved candidates produced zero passes and explored only four
families, converging on the same $\mathrm{ExpFamily}$ optimum the exhaustive
tier had already flagged as the closest near-miss and then stalling about
$53\times$ short of the bar.

\section{The referee broke before the theorem did}\label{sec:referee}

The most instructive thing that happened in this chapter was a failure of ours,
not of the theorem.

Within about ten candidates the evolutionary search drove its reported fitness
from $0.27$ to $0.46$, entirely inside $\mathrm{ExpFamily}$, by pushing the mean
per-offset cost from $3.71$ down to $2.20$ nats. It looked like a genuine climb
toward falsification. It was not: the search had found an accounting hole in our
scorer. $\mathrm{ExpFamily}$ was credited with about $10\,416$ certifiable
offsets while being charged $\log\mathrm{ord}$ nats \emph{per offset} with the
bill levied only against the $867$ residual primes---so a design was credited
with bulk coverage costing $10\,416\times2.20\approx22\,900$ nats of design
entropy against a total budget of $461$. A $50\times$ overdraft.

The audit protocol caught it before any claim was made, and the fix
\emph{strengthened} the mathematics. Charging per-offset mechanisms for every
offset they claim, and capping coverage at the affordable fraction
$\log B/\mathrm{cost}$, produced three things we did not have before: the
design-entropy budget constraint; the realisation that the honest benchmark is
covering's average $0.041$ nats per offset rather than its marginal $1.01$; and
the supply-side purchase argument of Observation~\ref{obs:expfamily}, which
closes $\mathrm{ExpFamily}$ quantitatively rather than empirically.

\begin{remark}[Methodological finding]\label{rem:referee}
Adversarial search against a formal claim is most valuable as a
\emph{referee-hardening} device, and should be reported that way rather than as
a failed refutation. The properties that made this work are worth restating,
because they are cheap and general: candidates emit data, never scores; grading
runs in a separate trusted process; every quantity is re-derived from
primitives; and a ``pass'' is a hand-audited event rather than an automatic
conclusion. Under those conditions an adversary optimising against your
evaluator is a free audit of your evaluator.
\end{remark}

\section{Verdict}

\begin{table}[h]
\centering\small
\caption{Per-family closure status. $13\,661$ enumerated plus $41$ evolved
schema instances produced zero passes; each family closes for its own
reason.}\label{tab:closure}
\begin{tabular}{lll}
\toprule
Family & Closure & Kind\\
\midrule
$\mathrm{Congruence}$ & is the covering baseline; composite $d$ dominated (P1) & proved\\
$\mathrm{PolyImage}$ ($k\ge2$) & value set $\le Q^{1/k}$, i.e. $\le349$ offsets & proved, price-free\\
$\mathrm{Cyclotomic}$ & $\varphi(n)\ge2$, same $\sqrt Q$ ceiling & proved, price-free\\
$\mathrm{ExpFamily}$ & purchase plan $1\,151$ nats vs $461$; cap $419$ & proved, quantitative\\
$\mathrm{MultiPoly}$ ($v<d$) & cluster bound $k_{\max}\le13$ & first moment\\
$\mathrm{MultiPoly}$ ($v\ge d$) & isolation + orbit conservation; else extraction & draft rigour\\
$\mathrm{NormForm}$ & reduces to $\mathrm{MultiPoly}$ & proved\\
$\mathrm{Compose}$ (depth 2) & intersection empty at $10^{200}$ & proved\\
\bottomrule
\end{tabular}
\end{table}

What remains genuinely conjectural, stated plainly:
\begin{enumerate}
\item \textbf{(H1)}, Hardy--Littlewood independence---standard but unproved,
supported here by the $1.2\times10^{9}$-candidate calibration.
\item \textbf{(H2)}, the sub-modexp detection question of
Conjecture~\ref{conj:h2}. A positive answer breaks the \emph{cost floor}, not
the covering monopoly.
\item The $o(1)$ terms in Lemma~\ref{lem:entropy}, and the
value-equidistribution inputs to the cluster counting.
\item \textbf{Mechanisms outside $L$.} The language is concrete and deliberately
finite. A mechanism that certifies compositeness by something that is neither a
planted divisor identity nor per-instance divisor discovery is not excluded by
anything here. We know of none, and the two structural horns---planted implies
entropy accounting, extracted implies factoring---suggest why, but this is the
honest frontier of the result.
\end{enumerate}

The programme-level consequence is worth stating because it changes what should
be worked on. If covering is not merely the best known paradigm but provably the
only competitive one within $L$, then records are governed entirely by the two
levers of Chapter~\ref{ch:frontier}: the rounding gap, and throughput. That is
the subject of the final chapter.

\chapter{Outlook: three doors and a costed path}\label{ch:outlook}

The preceding chapters were written in the order the work was done, which means
they read as a narrowing. That narrowing is the result. This chapter states it
positively: after Chapter~\ref{ch:closure}, everything that could possibly
improve a Goldbach-desert record passes through one of exactly three doors, and
each door has a maximal form, a price, and a verdict.

\section{The closure argument}

Why only three? Because the three ways to change $e^{E}$ are to change the
mathematics that sets $E$, to change the rate at which candidates are examined,
or to leave the paradigm in which $E$ is defined. Chapter~\ref{ch:frontier}
showed the first is a rounding gap, Chapter~\ref{ch:engineering} measured the
second, and Chapter~\ref{ch:closure} closed the third within a concrete
language. Two auxiliary facts complete the argument, and both are worth stating
because they eliminate whole families of otherwise-attractive ideas:

\begin{itemize}
\item \emph{Divisor information is fully harvested.} Any congruence cleverness
applied to the search parameter $k$ merely re-labels sub-progressions that the
sieve already scores. There is no free information on the $k$ side.
\item \emph{Correlation is covering} (Proposition~\ref{prop:shared}), and
beyond divisors the P5 measurement caps all remaining structure at $0.1$ nats.
So ``luck-shaping'' schemes have no purchase.
\end{itemize}

\section{Door 1: the rounding gap (mathematics)}

The exponent is entirely an integrality gap (Corollary~\ref{cor:rounding}), so
better rounding is worth records. Measured prize: greedy's $|U|=867$ against an
annealed ensemble threshold near $758$, i.e.\ $0$ to $2.6$ nats---a factor of up
to $14$, compounding into \emph{every} future rung.

Verdict from Chapter~\ref{ch:frontier}: resistant. Exact windows (LP-certified),
half-cover exact blocks, and BP decimation all leave $867$ untouched, and the
landscape is a field of isolated basins in which the annealed threshold is
probably not achievable. What has \emph{not} been tried is a different
\emph{ensemble} rather than a better search of the same one. Specifically:

\begin{enumerate}
\item \textbf{Banded construction with an exact endgame.} Solve the small moduli
exactly, choose the middle band with second-moment control on pairwise overlaps,
and finish the large moduli as an exact weighted matching against the
currently-uncovered set. This imitates \emph{why} greedy wins---sequential
adaptivity---while removing its myopia at the point where overlap waste is
greatest.
\item \textbf{Import Erd\H{o}s--Rankin technology constructively.} The
FGKMT machinery~\cite{FGKMT} is, at bottom, a method for rounding a fractional
cover of an interval with quantified loss; we used it in
Theorem~\ref{thm:lower} for a lower bound but never as a construction. Its loss
is $o(1)$ in an asymptotic regime that is not ours (their moduli reach $y/2$;
ours stop at $457$), so the question ``what residual fraction is achievable at
$\pi(Q)\approx10^{4}$ with a budget of $\theta(457)$?'' is genuinely open and
thesis-sized.
\item \textbf{One-step replica-symmetry-breaking.} Needed to get an honest
quenched threshold; naive belief propagation demonstrably slides toward the
fractional fantasy (\S\ref{sec:door1}).
\item \textbf{Post the instance as a public optimisation challenge.} The
benchmark is small, self-contained, and machine-checkable. Problems of exactly
this shape attract strong heuristic solvers when posed publicly.
\end{enumerate}

\section{Door 2: throughput (compute)}

Records are logarithmic in compute: each $10\times$ buys $2.3$ nats, which is
about four digits on the threshold game. Door 2 is therefore unglamorous and
completely reliable. Three forms, priced.

\paragraph{Engineering already in hand.} Porting the selective-testing scheme of
\S\ref{sec:engine} to the GPU fleet is worth $2$--$3\times$; batch
remainder-tree sieving~\cite{Bernstein} to depth $10^{9}$--$10^{10}$ is worth
perhaps $1.5\times$; replacing the fixed skip fraction with the exact marginal
stopping rule adds $5$--$15\%$. None of this is research.

\paragraph{Volunteer computing.} The search is a progress-free Poisson lottery
with millisecond verification and no shared state---structurally the best-shaped
volunteer workload that exists, and the exact shape on which
PrimeGrid~\cite{PrimeGrid} and its relatives hold essentially every record in
the neighbouring Sierpi\'nski--Riesel problems, which are themselves covering
problems. For calibration: a peak-era Folding@home-scale network is roughly
$2\times10^{5}$ times our two-GPU fleet, which would put $Q=200\,000$
(Table~\ref{tab:ladder}) about \emph{two days} away; even $1\%$ of that scale
delivers it in months.

\paragraph{Silicon.} The workload is $95\%$ one instruction---a
$660$-bit Montgomery modular exponentiation with a fixed base. A
crypto-mining-style ASIC plausibly buys $10$--$100\times$ an A100 per die at
much better energy per test, so a single rack is worth $e^{5.5}$ to $e^{8}$ over
today's fleet; an FPGA farm is the no-NRE middle step. This is the
capital-for-nats trade, and it is priced in the single-digit millions.

\section{Door 3: the paradigm}

Closed, within the language $L$ (Chapter~\ref{ch:closure}). The deliverable here
has therefore flipped from \emph{find a mechanism} to \emph{finish the theorem
and keep trying to break it}: complete the rigour of Lemma~\ref{lem:entropy},
settle (H2) or publicise it, and keep the falsification engine running as new
schema families are proposed. If the closure claim is ever refuted, the payoff
is enormous and immediate---but the sensible planning assumption is that it will
not be.

\section{Do we need a quantum computer?}

Grover search~\cite{Grover} halves the exponent: $e^{E}$ classical candidates
become $e^{E/2}$ coherent oracle calls. Since $E=37.7$ at $Q=200\,000$ and
$44.1$ for sub-$100$, this looks decisive. It is not, and the arithmetic is
worth doing once so that it need not be revisited.

One oracle call must evaluate the whole predicate reversibly: roughly $|U|$
modular exponentiations of $660$-bit numbers with \emph{no early exit}, since
the classical trick of stopping at the first passing complement is unavailable
in superposition. At fault-tolerant resource-estimate
scales~\cite{GidneyEkera} this is on the order of $10^{5}$--$10^{6}$ seconds per
call. A classical candidate costs $1.25\times10^{-6}$ seconds on our fleet.
Writing $C_q,W_q$ and $C_c,W_c$ for the per-operation costs and the widths
(numbers of machines), quantum wins only when
\[
E\;>\;E^{*}\;=\;2\log\Bigl[\frac{C_q}{W_q}\cdot\frac{W_c}{C_c}\Bigr].
\]

\begin{table}[h]
\centering\small
\caption{The Grover crossover. Compare against $E=37.7$ ($Q=200\,000$) and
$E=44.1$ (sub-$100$).}\label{tab:quantum}
\begin{tabular}{lr}
\toprule
Assumption & $E^{*}$\\
\midrule
One fault-tolerant machine, $C_q=10^{6}$\,s & $\approx55$\\
One thousand such machines & $\approx41$\\
One thousand machines \emph{and} a $1000\times$ cheaper oracle & $\approx27$\\
\bottomrule
\end{tabular}
\end{table}

With any non-miraculous constants the crossover sits at or above the hardest
targets we care about---and at the crossover both approaches are equally
infeasible. The honest answer is therefore: \emph{no}. You would need thousands
of much better quantum computers than anyone projects merely to \emph{tie} with
GPUs on this problem. Quantum walks for subset-sum face the same arithmetic in
Chapter~\ref{ch:barriers}'s setting. We park the idea with numbers attached
rather than dismissing it, so that it can be revisited if $C_q$ ever approaches
seconds.

\section{Operation Staircase: a costed path to \texorpdfstring{$\g>200\,000$}{g > 200000}}

Putting the doors together gives a concrete plan. Its shape is dictated by
Table~\ref{tab:ladder}: rungs are geometrically spaced, so climbing all of them
costs about $4\%$ more than jumping to the top, while banking a certified record
and a fresh calibration at every step.

\begin{enumerate}
\item \textbf{Now.} $Q=140\,009$ ($E\approx25$) on the existing fleet: days.
Bank it.
\item \textbf{Plus one week.} Land the engine work (selective testing on GPU,
marginal stopping rule): $2$--$3\times$. Then $Q=150\,000$ ($E=27.3$): about
four fleet-days. Bank it.
\item \textbf{Plus one month.} $Q\approx165\,000$ ($E\approx30$): two to three
fleet-weeks with the upgraded engine. In parallel, Door 1's banded construction
and the Erd\H{o}s--Rankin import run as theory work; any $\ge1$ nat result
compounds into every later rung.
\item \textbf{Plus one quarter.} $Q=175\,000$ ($E=32.1$): $1.4$ fleet-years raw,
so two to four months with the engine work and modest fleet growth. Decision
gate: if Door 1 has delivered $\ge2$ nats, stay in-house; otherwise open Door 2
at scale.
\item \textbf{The wall.} $Q=190\,000$--$200\,000$ ($E=35.4$--$37.7$) needs about
$e^{5.5}$ beyond step 4. Three interchangeable closers: a few hundred swarm or
volunteer GPUs for a quarter; an ASIC/FPGA rack; or four nats of Door-1
mathematics plus a $10\times$ fleet. Any one suffices; any two make it
comfortable.
\end{enumerate}

Expected profile: records at $140$k and $150$k within weeks, $165$k within a
month or two, $175$k inside a quarter, and $200\,000$ in two to four
quarters---with the long pole attackable by money, community, or mathematics,
whichever the programme can actually raise. Throughout, every rung is
dual-verified by the pipeline of \S\ref{sec:certification}, and parallel workers
claim disjoint variant sets in a committed ledger. That last point is not
bureaucratic: one of our own campaigns was stopped mid-flight for duplicating
another effort's target, which is a pure waste that coordination prevents.

\section{Open problems}

We close with the problems this thesis leaves in a state where someone else can
pick them up. The first two would change the subject; the rest would advance it.

\begin{enumerate}
\item \textbf{Sub-modexp compositeness detection} (Conjecture~\ref{conj:h2}).
Is there a randomised test with constant soundness on rough inputs, cheaper than
one modular exponentiation? Or a lower bound in any realistic model? A positive
answer speeds up every prime-hunting project in existence; a negative one
completes Lemma~\ref{lem:entropy}'s third leg.

\item \textbf{The rounding gap for prime-residue covers.} Given $\pi(Q)\approx
10^{4}$ prime offsets and moduli up to $457$, what residual fraction is
achievable? Greedy gives $7.56\%$, the fractional optimum $0\%$, and the annealed
ensemble threshold about $6.6\%$. Where is the quenched truth, and is it
constructive?

\item \textbf{CRT list-recovery beyond the Johnson radius} for structured,
solution-abundant instances (Chapter~\ref{ch:barriers}). Our experiments say the
radius is a sharp algorithmic boundary in practice; a proof or an algorithm
would both be significant.

\item \textbf{Modular subset-sum with per-position choices at density
$\approx1$}, where $2^{82}$ solutions exist among $2^{3893}$ designs.

\item \textbf{Mechanisms outside $L$}: certify $\N-q$ composite by something
that is neither a planted divisor identity nor per-instance divisor discovery.
This is the honest frontier of Chapter~\ref{ch:closure}'s claim.

\item \textbf{The smallest prescribed desert.} $S(q)$, the least even $\N$ with
$\g(\N)=q$~\cite{GLvdtR}, packages into $T(Q)=\min_{q>Q}S(q)$; no construction
here determines an exact minimum, and exhaustive verification stops at
$4\cdot10^{18}$ while our constructions start at $10^{150}$. The intervening
$130$ orders of magnitude are entirely unexplored.

\item \textbf{Closing the asymptotic sandwich.} Between
$\log\N\log_2\N\log_4\N/\log_3\N$ (rigorous, Theorem~\ref{thm:lower}) and
$(\log\N)^{2}\log\log\N$ (conjectural~\eqref{eq:GLR}) there is a factor of about
$\log\N$. Does covering only prime offsets---rather than a whole
interval---asymptotically beat the prime-gap scale?
\end{enumerate}

\section{A closing note on method}

This thesis was produced with heavy delegation to automated agents: search
campaigns, cover optimisation, literature checking, and an adversarial
falsification engine. Three practices did more than any individual result to
keep the output trustworthy, and we recommend them without reservation.

First, \emph{calibrate before believing}. The failure exponent was checked
against $1.2\times10^{9}$ candidates before it was used to make claims; the
belief-propagation implementation was validated against exhaustive enumeration
on a small instance before it was run on a large one. Both checks later caught
real errors.

Second, \emph{make negative results first-class}. Three campaigns produced no
record. Reported with their survival probabilities they became measurements
rather than embarrassments, and they now carry a substantial share of the
model's empirical support.

Third, \emph{let the adversary audit the referee}. The single most useful event
in Chapter~\ref{ch:closure} was an evolutionary search discovering a $50\times$
error in our own scorer. The audit protocol that caught it---candidates emit
data, never scores; grading in a separate trusted process; every quantity
re-derived from primitives; a ``pass'' is hand-audited---is cheap, general, and
worth more than the theorem it was built to defend. We also retracted two claims
in the course of this work (\S\ref{sec:retraction}, Remark~\ref{rem:units}) and
recorded both in place rather than quietly deleting them, which is the same
discipline pointed at ourselves.


\appendix
\chapter{Data, code, and reproduction}\label{app:data}

Every claim in this thesis that rests on computation is backed by an artifact in
the repository, and the artifacts are organised so that verification never
requires reproducing discovery. This appendix is the index.

\section{Certified records}

Each record directory contains the canonical decimal string of $\N$, the
congruence cover, the reconstruction data $(\N_0,M,k)$, per-offset exclusion
evidence, an Atkin--Morain ECPP certificate for the prime complement, and a
SHA-256 manifest.

\begin{center}\small
\begin{tabular}{lll}
\toprule
Directory & Record & Verification\\
\midrule
\texttt{records/megagap\_2692digit\_g1134871/} & $2\,692$d, $\g=1\,134\,871$ & witnesses + ECPP\\
\texttt{records/threshold\_150digit\_g104527/} & $150$d, $\g=104\,527$ & witnesses + ECPP, re-verified\\
\texttt{records/threshold\_195digit\_g100297/} & $195$d, $\g=100\,297$ & witnesses + ECPP\\
\texttt{records/dual\_197digit\_g107719/} & $197$d, $\g=107\,719$ & witnesses + ECPP\\
\texttt{records/budget\_1e200\_200digit\_g112249/} & $200$d, $\g=112\,249$ & witnesses + ECPP\\
\texttt{slop/goldbach/records/g1/} & $2\,480$d, $\g=1\,157\,341$ & witnesses + ECPP\\
\texttt{slop/goldbach/records/r200/} & $199$d, $\g=119\,419$ & witnesses + ECPP, re-verified\\
\texttt{slop/goldbach/records/t100k/} & ratchet $195\to179$d & witnesses + ECPP\\
\bottomrule
\end{tabular}
\end{center}

Two independent verification paths exist: \texttt{verify\_record.py}
re-derives every per-offset witness from the record file alone, and
\texttt{ecpp\_verify.py} validates the elliptic-curve certificate chain
step by step in projective coordinates without invoking the library that
generated it.

\section{Search engine and campaigns}

\begin{description}
\item[\texttt{goldbach/search.py}] The CRT-progression search engine: numpy
presieve with cached modular inverses, survivor-ordered testing
(\texttt{--sort-tests}), and selective testing with exact hit-mass accounting
(\texttt{--skip-frac}, \texttt{--p-est}). Reports per-slice $\hat p$,
survivors per candidate, $\Ehat$, and retained hit-mass share---the telemetry
behind \S\ref{sec:calibration}.
\item[\texttt{goldbach/cover.py}, \texttt{anneal.py}, \texttt{joint.py}]
Cover construction: greedy weighted set cover, simulated annealing, and the
joint representative/coverage optimiser.
\item[\texttt{goldbach/gpu/}] The CUDA engine (bit-matrix sieve, Montgomery CIOS
Fermat waves) and the multi-session fleet driver with git-committed resume
marks.
\item[\texttt{goldbach/phd/campaign\{,2,3\}/}] The three campaigns of
Table~\ref{tab:campaigns}: specs, chunked drivers, and the committed
\texttt{search.log} checkpoints. The logs are the primary evidence for the
negative results, and they make each campaign resumable.
\end{description}

\section{Frontier and optimisation studies}

\begin{description}
\item[\texttt{goldbach/phd/wp8/pack\_probe.py}] Coordinate ascent with random
restarts over class assignments; produces the rung ladder of
Table~\ref{tab:ladder} and the local-optimality finding.
\item[\texttt{goldbach/phd/wp8/lp\_bound.py}] The set-cover LP relaxation
(HiGHS) behind Proposition~\ref{prop:lp}.
\item[\texttt{goldbach/phd/wp9/door1/}] The Door-1 benchmark and attacks:
\texttt{instance\_q122k.json} (sha256-pinned prime universe, the $|U|=867$
incumbent), \texttt{d10\_diagnostics.py} (landscape statistics),
\texttt{d11\_lns.py} (CP-SAT large-neighbourhood search with per-window LP
certificates), \texttt{d11\_basin\_probe.py}, and \texttt{d12\_bp.py} (belief
propagation with decimation, including the exhaustive-enumeration validation).
\end{description}

\section{Barrier and closure studies}

\begin{description}
\item[\texttt{goldbach/phd/wp3\_small\_representative\_memo.md}] The sub-$100$
analysis of Chapter~\ref{ch:barriers}, including the retraction of
\S\ref{sec:retraction} and the corrected costing of Remark~\ref{rem:units}.
\item[\texttt{goldbach/phd/toy\_crt/}] The CRT-decoding testbed: planted
instances, LLL decoding, and the threshold sweep across the Johnson-type radius.
\item[\texttt{goldbach/phd/wp7\_probes.md}] The five mechanism probes of
\S\ref{sec:wp7} with their predeclared gates.
\item[\texttt{goldbach/phd/wp9/wp9\_closure\_theorem.md}] The closure theorem in
full working form, with proof-status markers throughout.
\item[\texttt{goldbach/phd/wp9/engine/}] The falsification engine:
\texttt{tier0.py} (calibration against the hand-derived numbers),
\texttt{exhaustive.py} (the $13\,661$-instance sweep) with
\texttt{exhaustive\_summary.json} (per-family aggregates and every near-miss;
the full rows are regenerable and deliberately not committed),
\texttt{scorer.py} (the trusted evaluator, with the budget accounting of
\S\ref{sec:referee}), \texttt{harness.md} (the anti-reward-hack protocol), and
\texttt{ae\_runner.py} with \texttt{ae\_seed\_program.py} (the generative tier).
\item[\texttt{goldbach/phd/wp9/h2\_open\_problem.md}] Conjecture~\ref{conj:h2},
stated self-contained for posting.
\end{description}

\section{Reproducing the numerical claims}

The claims of Chapters~\ref{ch:frontier}--\ref{ch:closure} are reproducible from
the committed scripts, and all are cheap---seconds to minutes---because they are
analyses rather than searches:

\begin{center}\small
\begin{tabular}{ll}
\toprule
Claim & Command\\
\midrule
LP vacuousness (Prop.~\ref{prop:lp}) & \texttt{python3 phd/wp8/lp\_bound.py 122000 200000}\\
Rung ladder (Table~\ref{tab:ladder}) & \texttt{python3 phd/wp8/pack\_probe.py}\\
Door-1 diagnostics & \texttt{python3 phd/wp9/door1/d10\_diagnostics.py}\\
Exact-window search & \texttt{python3 phd/wp9/door1/d11\_lns.py 15 7 3 30}\\
BP validation + decimation & \texttt{python3 phd/wp9/door1/d12\_bp.py}\\
Engine calibration (tier 0) & \texttt{python3 phd/wp9/engine/tier0.py}\\
Exhaustive schema sweep & \texttt{python3 phd/wp9/engine/exhaustive.py}\\
Single-schema score & \texttt{python3 phd/wp9/engine/scorer.py '\{...\}'}\\
\bottomrule
\end{tabular}
\end{center}

Dependencies are \texttt{numpy}, \texttt{sympy}, \texttt{gmpy2},
\texttt{scipy} (HiGHS), and \texttt{ortools} (CP-SAT); the record verifiers
require only Python 3, deliberately, so that verification has no dependency on
the discovery stack.

\section{A note on provenance}

The programme ran as several parallel efforts sharing one repository, and the
records in Table~\ref{tab:records} were not all produced by the same effort. We
have tried to be accurate about this throughout: where a record was produced
elsewhere in the programme and re-verified here, the text says so, and the
research log (\texttt{goldbach/phd/RESEARCH\_LOG.md}) records the day-by-day
provenance including the campaigns that failed, the two retractions, and the
coordination failure that stopped Campaign~3. The theoretical contributions of
Chapters~\ref{ch:frontier}--\ref{ch:closure}---the LP vacuousness result, the
Door-1 optimisation study, the sub-$100$ barrier analysis with its retraction,
and the closure theorem with its falsification engine---were produced by the
work reported here.


\backmatter
\begin{thebibliography}{99}
\bibitem{OeS} T.~Oliveira e Silva, S.~Herzog, S.~Pardi, \emph{Empirical
verification of the even Goldbach conjecture and computation of prime gaps up to
$4\cdot10^{18}$}, Math. Comp. \textbf{83} (2014), 2033--2060.
\bibitem{Richstein} J.~Richstein, \emph{Verifying the Goldbach conjecture up to
$4\cdot10^{14}$}, Math. Comp. \textbf{70} (2001), 1745--1749.
\bibitem{HL} G.~H.~Hardy, J.~E.~Littlewood, \emph{Some problems of `Partitio
Numerorum' III}, Acta Math. \textbf{44} (1923), 1--70.
\bibitem{GLvdtR} A.~Granville, J.~van de Lune, H.~J.~J.~te Riele, \emph{Checking
the Goldbach conjecture on a vector computer}, in: Number Theory and Applications
(R.~A.~Mollin, ed.), NATO ASI Ser.~C \textbf{265}, Kluwer, Dordrecht, 1989,
423--433.
\bibitem{OEIS} N.~J.~A.~Sloane (ed.), \emph{The On-Line Encyclopedia of Integer
Sequences}, sequences A020481 (least Goldbach summand) and A025018/A025019
(records), \url{https://oeis.org}.
\bibitem{JB} Zs.~Juh\'asz, M.~Bartalos, \emph{Predicting the size ranking of
minimal primes in the generalised Goldbach partitions}, arXiv:2510.21870 (2025).
\bibitem{Dirichlet} P.~G.~L.~Dirichlet, \emph{Beweis des Satzes, dass jede
unbegrenzte arithmetische Progression\dots\ unendlich viele Primzahlen
enth\"alt}, Abh. K. Preuss. Akad. Wiss. (1837), 45--81.
\bibitem{Linnik} U.~V.~Linnik, \emph{On the least prime in an arithmetic
progression I, II}, Mat. Sbornik \textbf{15(57)} (1944), 139--178, 347--368.
\bibitem{Rankin} R.~A.~Rankin, \emph{The difference between consecutive prime
numbers}, J. London Math. Soc. \textbf{13} (1938), 242--247.
\bibitem{Erdos1950} P.~Erd\H{o}s, \emph{On integers of the form $2^k+p$ and
some related problems}, Summa Brasil. Math. \textbf{2} (1950), 113--123.
\bibitem{FGKT} K.~Ford, B.~Green, S.~Konyagin, T.~Tao, \emph{Large gaps between
consecutive prime numbers}, Ann. of Math. \textbf{183} (2016), 935--974.
\bibitem{Maynard} J.~Maynard, \emph{Large gaps between primes}, Ann. of Math.
\textbf{183} (2016), 915--933.
\bibitem{FGKMT} K.~Ford, B.~Green, S.~Konyagin, J.~Maynard, T.~Tao, \emph{Long
gaps between primes}, J. Amer. Math. Soc. \textbf{31} (2018), 65--105.
\bibitem{PSW} C.~Pomerance, J.~L.~Selfridge, S.~S.~Wagstaff~Jr., \emph{The
pseudoprimes to $25\cdot10^9$}, Math. Comp. \textbf{35} (1980), 1003--1026.
\bibitem{BW} R.~Baillie, S.~S.~Wagstaff~Jr., \emph{Lucas pseudoprimes},
Math. Comp. \textbf{35} (1980), 1391--1417.
\bibitem{GoldwasserKilian} S.~Goldwasser, J.~Kilian, \emph{Almost all
primes can be quickly certified}, Proc. 18th ACM Symposium on Theory of
Computing (1986), 316--329.
\bibitem{AtkinMorain} A.~O.~L.~Atkin, F.~Morain, \emph{Elliptic curves and
primality proving}, Math. Comp. \textbf{61} (1993), 29--68.
\bibitem{PARI} The PARI Group, \emph{PARI/GP version 2.17.2}, 2026,
\url{https://pari.math.u-bordeaux.fr}.
\bibitem{GSS} V.~Guruswami, A.~Sahai, M.~Sudan, \emph{``Soft-decision'' decoding
of Chinese Remainder codes}, Proc. 41st IEEE Symposium on Foundations of
Computer Science (2000), 159--168.
\bibitem{Goldreich} O.~Goldreich, D.~Ron, M.~Sudan, \emph{Chinese remaindering
with errors}, IEEE Trans. Inform. Theory \textbf{46} (2000), 1330--1338.
\bibitem{LagariasOdlyzko} J.~C.~Lagarias, A.~M.~Odlyzko, \emph{Solving
low-density subset sum problems}, J. ACM \textbf{32} (1985), 229--246.
\bibitem{Coster} M.~J.~Coster, A.~Joux, B.~A.~LaMacchia, A.~M.~Odlyzko,
C.~P.~Schnorr, J.~Stern, \emph{Improved low-density subset sum algorithms},
Comput. Complexity \textbf{2} (1992), 111--128.
\bibitem{Wagner} D.~Wagner, \emph{A generalized birthday problem}, Advances in
Cryptology---CRYPTO 2002, LNCS \textbf{2442}, 288--304.
\bibitem{LLL} A.~K.~Lenstra, H.~W.~Lenstra~Jr., L.~Lov\'asz, \emph{Factoring
polynomials with rational coefficients}, Math. Ann. \textbf{261} (1982),
515--534.
\bibitem{Grover} L.~K.~Grover, \emph{A fast quantum mechanical algorithm for
database search}, Proc. 28th ACM Symposium on Theory of Computing (1996),
212--219.
\bibitem{GidneyEkera} C.~Gidney, M.~Eker\aa, \emph{How to factor 2048 bit RSA
integers in 8 hours using 20 million noisy qubits}, Quantum \textbf{5} (2021),
433.
\bibitem{Bernstein} D.~J.~Bernstein, \emph{How to find smooth parts of
integers}, preprint, 2004, \url{https://cr.yp.to/factorization/smoothparts.pdf}.
\bibitem{MoserTardos} R.~A.~Moser, G.~Tardos, \emph{A constructive proof of the
general Lov\'asz Local Lemma}, J. ACM \textbf{57} (2010), article 11.
\bibitem{MezardMontanari} M.~M\'ezard, A.~Montanari, \emph{Information, Physics,
and Computation}, Oxford University Press, 2009.
\bibitem{CPSAT} L.~Perron, F.~Didier, \emph{CP-SAT}, Google OR-Tools, 2024,
\url{https://developers.google.com/optimization/cp/cp_solver}.
\bibitem{Cornacchia} G.~Cornacchia, \emph{Su di un metodo per la risoluzione in
numeri interi dell'equazione $\sum_{h=0}^{n}C_hx^{n-h}y^h=P$}, Giornale di
Matematiche di Battaglini \textbf{46} (1908), 33--90.
\bibitem{Zsygmondy} K.~Zsigmondy, \emph{Zur Theorie der Potenzreste},
Monatshefte f\"ur Mathematik \textbf{3} (1892), 265--284.
\bibitem{AlphaEvolve} Google DeepMind, \emph{AlphaEvolve: a coding agent for
algorithmic discovery}, developer documentation, 2026,
\url{https://docs.cloud.google.com/gemini/enterprise/docs/alphaevolve}.
\bibitem{Sierpinski} W.~Sierpi\'nski, \emph{Sur un probl\`eme concernant les
nombres $k\cdot2^n+1$}, Elem. Math. \textbf{15} (1960), 73--74.
\bibitem{PrimeGrid} PrimeGrid, \emph{Distributed computing project},
\url{https://www.primegrid.com}.
\end{thebibliography}

\end{document}
